% Standalone source generated from article.tex by build.py.
% All mathematical statements, proofs, and bibliography are included.

% BEGIN preamble.tex
\documentclass[11pt,a4paper]{article}
\usepackage[T1]{fontenc}
\usepackage[utf8]{inputenc}
\usepackage{lmodern}
\usepackage[margin=25mm,headheight=15pt]{geometry}
\usepackage{amsmath,amssymb,amsthm,mathtools}
\usepackage{booktabs,array,longtable,tabularx}
\usepackage{microtype}
\usepackage{xcolor}
\usepackage{enumitem}
\usepackage{needspace}
\usepackage{fancyhdr}
\usepackage{xurl}
\usepackage[colorlinks=true,linkcolor=blue!45!black,citecolor=blue!45!black,
            urlcolor=blue!50!black,bookmarksnumbered=true]{hyperref}
\usepackage{bookmark}
\numberwithin{equation}{section}
\newtheorem{theorem}{Theorem}[section]
\newtheorem{proposition}[theorem]{Proposition}
\newtheorem{lemma}[theorem]{Lemma}
\newtheorem{corollary}[theorem]{Corollary}
\theoremstyle{definition}
\newtheorem{definition}[theorem]{Definition}
\newtheorem{conjecture}[theorem]{Conjecture}
\theoremstyle{remark}
\newtheorem{remark}[theorem]{Remark}
\DeclareMathOperator{\Li}{Li}
\DeclareMathOperator{\Log}{Log}
\newcommand{\src}[1]{\texttt{\detokenize{#1}}}
\newcommand{\snapshot}{\texttt{c2cb285b64273e4cb2c8d7a03c68a4a99ac812d4}}
\newcolumntype{Y}{>{\raggedright\arraybackslash}X}
\setlist{itemsep=3pt,topsep=6pt}
\setlength{\emergencystretch}{1.5em}
\setlength{\parskip}{3pt}
\allowdisplaybreaks[2]
\pagestyle{fancy}
\fancyhf{}
\fancyhead[L]{\small Harmonic parity and resolvent identities}
\fancyhead[R]{\small ProveIt research continuation}
\fancyfoot[C]{\thepage}
\renewcommand{\headrulewidth}{0.3pt}
\hypersetup{
 pdftitle={Harmonic Parity and Resolvent Identities},
 pdfauthor={Research continuation for ProveIt},
 pdfsubject={Exact harmonic reductions, polylogarithmic transforms, Stieltjes primitives, and Tornheim jets},
 pdfkeywords={polylogarithm, Hurwitz zeta, harmonic numbers, Stieltjes constants, polygamma, Tornheim}
}

% END preamble.tex

\begin{document}
\begin{titlepage}
\thispagestyle{empty}
\vspace*{12mm}
\begin{center}
{\small\scshape Exact identities and analytic continuation\par}
\vspace{9mm}
{\LARGE\bfseries Harmonic Parity and\\[3pt] Resolvent Identities\par}
\vspace{6mm}
{\large Exact Lerch generators, Stieltjes primitives,\\
and the quartic and higher Tornheim layers\par}
\vspace{10mm}
{\normalsize Research continuation for the ProveIt project\par}
\vspace{3mm}
{\normalsize 11 October 2026\par}
\end{center}
\vspace{8mm}

% BEGIN sections/00_abstract.tex
\begin{abstract}
We develop four exact continuations of the ProveIt program on
polylogarithms, harmonic numbers, Gamma functions, and zeta derivatives.
First, for an arbitrary polynomial numerator in generalized harmonic
numbers, we classify the cancellation of every nonpositive even
spectral principal part at the centered shift $1/2$. The criterion is
invariance under the simultaneous reflection of all even harmonic
tails; its necessity follows from a theorem on formal difference
equations. We then evaluate every centered even moment of the square
of trigamma as a rational number plus a rational multiple of $\zeta(3)$,
and sum these values in a convergent polylogarithmic generator.

Second, a cotangent resolvent gives a completed formula for all
generalized Stieltjes indices and all argument derivatives. We
identify the entire local correction polynomial before coefficient
extraction. Partial fractions yield every rational cotangent kernel
bounded at infinity with no real pole. The same transform handles
continuous polygamma orders and all balanced Gamma primitives, with
an explicit conversion between analytic continuation in the order
and pointwise endpoint finite parts.

Third, fixed nonpositive indices can be eliminated at arbitrary
positions and depths in the entire relative harmonic interpolant.
We prove a canonical finite polynomial normal form and a genuinely
convergent Lerch generator for arbitrarily many zero indices surrounding
one free order. Classical polynomial Stieltjes primitives are recorded
with arbitrary endpoints and fully explicit coefficients.

Finally, every admissible quartic Tornheim ray is reduced to three
specified coordinates, and cyclic quartic combinations retain only
the diagonal coordinate. At all degrees we determine the exact
homogeneous freedom under the stated symmetry, axis, and Euler-plane
relations. An explicit deformation proves why diagonal data cease
to determine the cyclic fifth layer. The article includes complete
proofs, a source audit and integration notes, reproducible exact and numerical
checks, and a concrete further-research agenda.
\end{abstract}

% END sections/00_abstract.tex

\vfill
\noindent\textbf{Source baseline.}
The manuscript and incoming reports are inspected at the fixed repository
revision \snapshot. Proof status and contributions are stated relative to
that baseline. This article does not claim a proof of the unresolved
$S_6$ or revised $S_8$ reductions.
\end{titlepage}
\tableofcontents
\clearpage

% BEGIN sections/01_scope.tex
\section{Scope, main results, and mathematical conventions}
\label{scope:section}

The objective is to obtain exact identities that survive the passage
between spectral parameters, endpoint singularities, and ordinary
special functions. Four aspects of the existing research collection
are particularly productive: centered harmonic expansions,
cotangent--Hurwitz transforms, entire relative harmonic interpolants,
and holomorphic remainders of multivariate zeta functions. Each already
has a substantial foundation in the manuscript or the incoming reports.
The present article extends that foundation at specific open points.

\subsection{The inspected baseline and the status of the questions}

The source baseline is the repository revision \snapshot,
dated 11 October 2026 at 00:52:39 UTC. We inspected the relevant
canonical manuscript chapters and the nine incoming archives listed
in the accompanying source manifest. The most direct starting points
are \emph{Exact Resonant Identities} \cite{ghp:ExactResonant},
\emph{Independent Orders and Exact Reductions} \cite{qtr:Independent},
and \emph{Cubic Tornheim Derivatives and Harmonic Stieltjes
Constants} \cite{source:Cubic}. The original manuscript remains
the reference for its established normalization and conjecture labels
\cite{source:manuscript}.

This is a targeted continuation and audit of those areas, not a claim
to have checked every line of all manuscript volumes. In particular,
the short $S_6$ reduction and the revised $S_8$ reduction remain
conjectures. Previously rejected candidate vectors and previously
reported corrections are distinguished from current candidates in
Section~\ref{audit:section}.

The following table states the exact scope of the main advances.
Resolution refers to the indicated mathematical question, with the
domains and qualifications in the cited theorem.

\begin{table}[htbp]
\centering
\small
\begin{tabularx}{\textwidth}{@{}p{0.23\textwidth}Y p{0.19\textwidth}@{}}
\toprule
Source question or direction & Result proved here & Location\\
\midrule
Exact Resonant Q11: generalized harmonic cancellation
& Necessary and sufficient centered involution criterion for all
nonpositive even spectral points.
& Theorem~\ref{ghp:thm:classification}\\
\addlinespace
Explicit centered values beyond raw $H_n$ powers
& Every trigamma-square even moment, its ordinary convergent
subtracted sum, and a convergent polylogarithmic generator.
& Theorems~\ref{ghp:thm:square}, \ref{ghp:thm:generator}\\
\addlinespace
Exact Resonant Q8: a nonpolynomial cotangent family
& The complete rational sector bounded at infinity, with the
all-index endpoint correction.
& Theorem~\ref{res:main}; Corollary~\ref{res:repeated}\\
\addlinespace
Exact Resonant Q9: normalized primitive families
& The resolvent transform of every balanced Gamma primitive and
the continuous-order conversion formula.
& Theorem~\ref{res:continuous}; Corollary~\ref{res:balancedcor}\\
\addlinespace
Independent Orders: nonpositive-index reductions
& Arbitrary positions and depths, a canonical finite normal form,
and an analytic zero-decoration generator.
& Theorems~\ref{harm:local}--\ref{harm:Lerch}\\
\addlinespace
Independent Orders and Cubic Harmonic: quartic rays
& All quartic directions in three coordinates; cyclic directions
in one diagonal coordinate.
& Theorem~\ref{qtr:quartic}; equation~\eqref{qtr:cyclicquartic}\\
\addlinespace
Higher directional relations
& Exact formal kernel and cyclic-image dimensions for the named
relations, and a fifth-order obstruction.
& Theorem~\ref{qtr:kernel}; Corollary~\ref{qtr:cyclicimage}\\
\bottomrule
\end{tabularx}
\caption{Results relative to the pinned source questions. The primitive
and higher-order statements resolve the specified families, rather than
every possible extension proposed in the source.}
\label{scope:table}
\end{table}

\subsection{A selection of exact identities}

Three compact examples illustrate the different kinds of conclusion.
The centered square identity
\begin{equation}\label{scope:sample-square}
 \sum_{n=0}^{\infty}
 \left((n+\tfrac12)^2\psi'(n+1)^2-1\right)
       =-\frac14\zeta(3)-\frac16
\end{equation}
is an ordinary absolutely convergent series. It is the first
nontrivial member of a complete all-even family, rather than a
numerically recognized isolated value.

The pointwise endpoint finite part
\begin{equation}\label{scope:sample-resolvent}
 \operatorname{FP}_x\int_0^1
 \frac{\psi'(x)}{\pi\cot(\pi x)-i\pi}\,dx
       =1-\gamma-\log(2\pi)+\frac{i\pi}{2}
\end{equation}
has a boundary contribution equal to $1$. Analytically continuing an
integrated spectral family and then substituting the integer order
without that contribution gives a different value. The all-index
theorem locates this difference before any expansion.

For the Tornheim restriction $W_{a,b,c}(t)=T(at,bt,ct)$ and
$\Lambda=W_{1,1,1}^{(4)}(0)$, a cyclic combination gives
\begin{align}\label{scope:sample-quartic}
 &W_{1,2,3}^{(4)}(0)+W_{2,3,1}^{(4)}(0)
       +W_{3,1,2}^{(4)}(0)-36\Lambda\notag\\
 &\hspace{13mm}=-184\log(2\pi)\zeta'''(0)
                +156\zeta''(0)^2-386\zeta''''(0).
\end{align}
The cancellation is an exact polynomial identity between directional
coefficients. It does not require an assumption about relations among
the individual constants.

\subsection{Conventions and the meaning of a reduction}

We use
\[
 \zeta(1+t,x)=\frac1t+
      \sum_{m\ge0}\frac{(-1)^m}{m!}\gamma_m(x)t^m,\qquad
 \gamma_m=\gamma_m(1),\qquad \gamma=\gamma_0,
\]
and $\psi=\Gamma'/\Gamma$. Primes on a one-variable zeta function, or
$\zeta^{(j)}(s,x)$ with the first argument displayed, mean derivatives
in the spectral argument. Superscripts on $\psi$ and $\gamma_m(x)$
mean derivatives in $x$. When confusion is possible, a square
superscript $[j]$ denotes a polylogarithmic order derivative:
\[
 \Li_s^{[j]}(z)=\partial_s^j\Li_s(z).
\]
All logarithms of positive real quantities are real. Complex powers
use the specified branch; Hurwitz endpoints are initially in the
right half-plane. Bernoulli numbers satisfy $B_1=-1/2$.
All branches required by the resolvent formulas are given explicitly
in Section~\ref{res:section}.

The term \emph{reduction} is always relative to a stated target class.
For the harmonic deletion theorem it means a finite linear combination
with polynomial endpoint coefficients and only the surviving free
orders. For a rational cotangent kernel it means finitely many
polylogarithmic order derivatives at explicit arguments. For the
quartic Tornheim layer it means three explicitly normalized
coordinates, two of which are supplied with ordinary convergent
representatives. It does not silently mean reduction to a conjectured
algebra of familiar constants.

There are three distinct analytic operations in the proofs.
A meromorphic continuation is uniquely determined from a convergence
domain. A coordinate finite part extracts the constant term of a
cutoff integral in its stated endpoint variable. A directional
restriction of a meromorphic function is formed before taking its
one-variable limit. These operations sometimes agree and sometimes
differ by explicit terms. The article proves the needed compatibility
at each use.

\subsection{Proof inputs and attribution}

The sufficient half-shift parity argument is elementary asymptotic
algebra. Its converse uses the formal-Laurent version of the
Adamczewski--Dreyfus--Hardouin theorem
\cite[Theorem~1.2]{ghp:ADH}. This use is spelled out, including a
proof that the relevant formal trigamma series is not rational.
It supplies differential algebraic independence of functions and
formal series; it says nothing about algebraic independence of
Euler's constant or zeta values.

The Mellin, Fourier, Bernoulli, Gamma, and polylogarithm expansions
used below are classical \cite{DLMF}. The relative harmonic
interpolant and the Tornheim polar normalization are inherited from
the pinned incoming reports; the underlying one-endpoint
interpolant is due to Ihara, Nakamura, and Yamamoto \cite{harm:INY}.
The continuous polygamma function and balanced negapolygamma functions
are the established Espinosa--Moll functions \cite{EM}.
Polynomial Hurwitz and Stieltjes primitives are also established
in Espinosa--Moll \cite{harm:EM} and the canonical manuscript.
Our fully indexed arbitrary-endpoint version is included to make the
new finite reduction calculus directly usable.

The centered trigamma-square family should also be distinguished
from the weighted squares of \emph{differences} of trigamma functions
already evaluated in the incoming mixed-spectral report
\cite{source:Mixed}. Both are exact identities; their summands,
parameters, and generating mechanisms are different.

All results described as developments are developments relative to
the inspected source baseline. The article makes no literature-wide
priority claim. Proofs are analytic and algebraic; the accompanying
programs provide finite exact checks and numerical diagnostics, not
proof-assistant verification of the analytic theorems.

% END sections/01_scope.tex


% BEGIN sections/02_harmonic_parity.tex
% Manuscript-ready section. Requires amsmath, amssymb, amsthm.
% Bibliographic keys: ghp:ADH, ghp:DLMF, ghp:ExactResonant.
\section{Exact parity of generalized harmonic numerators}
\label{ghp:sec}

Question Q11 in the pinned \emph{Exact Resonant Identities} report
\cite{ghp:ExactResonant} asks which polynomials in generalized harmonic
numbers have no principal part at any nonpositive even spectral integer
when the outer shift is $1/2$. The raw-power argument proves a sufficient
condition for powers of $H_n$; it does not itself classify arbitrary
polynomials. This section gives a necessary and sufficient algebraic
criterion. The necessity uses an established theorem on formal
difference equations. It does not assume algebraic independence of
Euler's constant, zeta values, or numerical periods.

\subsection{The centered involution and the classification theorem}

Put $H_0^{(r)}=0$, $H_n^{(r)}=\sum_{k=1}^n k^{-r}$, and $H_n=H_n^{(1)}$.
For a fixed $R\geq1$ and a polynomial $P\in\mathbb C[X_1,\ldots,X_R]$,
let
\begin{equation}\label{ghp:def:series}
 \mathcal D_P(s,a)=\sum_{n=0}^{\infty}
  \frac{P(H_n,H_n^{(2)},\ldots,H_n^{(R)})}{(n+a)^s},
 \qquad a>0,\quad \Re s>1.
\end{equation}
These series converge normally in that half-plane; a polynomial
numerator grows at most as a fixed power of $\log(n+2)$.
All later values refer to the unique meromorphic continuation.

Define an involution on the polynomial ring by
\begin{equation}\label{ghp:def:tau}
 \tau X_1=X_1,\qquad
 \tau X_r=\begin{cases}
       2\zeta(r)-X_r,&r\geq2\text{ even},\\
       X_r,&r\geq3\text{ odd}.
       \end{cases}
\end{equation}
The coefficients in $\mathbb C$ are fixed. Thus the transformation
changes the sign of every even-order tail $\zeta(2j)-X_{2j}$
\emph{simultaneously}.

\Needspace{20\baselineskip}
\begin{theorem}[Exact centered cancellation criterion]\label{ghp:thm:classification}
For $P\in\mathbb C[X_1,\ldots,X_R]$, the following conditions are
 equivalent:
\begin{enumerate}
\item $\mathcal D_P(s,1/2)$ is holomorphic at every
      $s=0,-2,-4,\ldots$.
\item $\tau P=P$ as a polynomial with complex coefficients.
\item On writing $V_j=\zeta(2j)-X_{2j}$, every monomial of $P$ has
      even total degree in the variables $V_1,V_2,\ldots$.
\end{enumerate}
Equivalently, the complete algebra of such numerators is
\begin{equation}\label{ghp:eq:invariant-ring}
 \mathbb C\bigl[X_1,X_3,X_5,\ldots,
  (\zeta(2i)-X_{2i})(\zeta(2j)-X_{2j}):1\leq i\leq j\leq\lfloor R/2\rfloor\bigr].
\end{equation}
Only variables whose index is at most $R$ are retained.
\end{theorem}

The paired generators in \eqref{ghp:eq:invariant-ring} are a generating
set, not an assertion that this algebra is freely generated. For
example, $(V_iV_j)^2=V_i^2V_j^2$ is an algebraic relation among them.

Examples of admissible numerators are
\[
 H_n^pH_n^{(3)},\qquad
 (\zeta(2)-H_n^{(2)})^2,\qquad
 H_n^p(\zeta(2)-H_n^{(2)})(\zeta(4)-H_n^{(4)}),
 \qquad p\geq0.
\]
The raw polynomial $H_n^{(2)}$ does not qualify. In fact its principal
part at $s=0$ is $-1/s$, since
$H_n^{(2)}=\zeta(2)-(n+1/2)^{-1}+O(n^{-3})$.
Multiplying two even-order tails changes this outcome by making their
full centered expansion even.

\subsection{A finite formula for every principal part}

Let $L$ be an independent formal symbol, $z$ a formal variable, and
\begin{equation}\label{ghp:def:A}
 A(z,a)=-\sum_{k\geq1}\frac{B_k(a)}kz^k,
 \qquad T_1(z,L;a)=L+\gamma+A(z,a).
\end{equation}
For $r\geq2$ define
\begin{equation}\label{ghp:def:Tr}
 T_r(z;a)=\zeta(r)-\frac{z^{r-1}}{r-1}
 -\frac1{r-1}\sum_{k\geq1}
   \binom{k+r-2}{r-2}B_k(a)z^{k+r-1}.
\end{equation}
These are formal asymptotic series; convergence in $z$ is not needed.
Set
\begin{equation}\label{ghp:def:q}
 P(T_1,T_2,\ldots,T_R)=\sum_{j\geq0}q_{P,j}(L;a)z^j,
 \qquad q_{P,j}(L;a)=\sum_{\ell\geq0}q_{P,j,\ell}(a)L^\ell.
\end{equation}
Every coefficient requires only finitely many Bernoulli polynomials and
finite polynomial operations.

\begin{proposition}[Generalized-harmonic principal coefficients]
\label{ghp:prop:principal}
The function \eqref{ghp:def:series} continues meromorphically to the
entire $s$-plane. Its possible poles are $1,0,-1,-2,\ldots$. At
$s=1-j$ its complete principal part is
\begin{equation}\label{ghp:eq:principal}
 \sum_{\ell\geq0}\frac{\ell!\,q_{P,j,\ell}(a)}{(s+j-1)^{\ell+1}}.
\end{equation}
The sum is finite. In particular, this point is regular if and only if
$q_{P,j}(L;a)$ is the zero polynomial in $L$.
\end{proposition}

\begin{proof}
Writing $x=n+a$, the ordinary shifted digamma expansion and its fixed
argument derivatives give the expansions $T_r$ above, with $z=x^{-1}$
and $L=\log x$; see the classical Gamma and polygamma formulas in
\cite[\S\S5.11, 5.15]{DLMF}. For $r\geq2$ one can also obtain
\eqref{ghp:def:Tr} directly by differentiating
\[
 H_n=\psi(n+1)+\gamma,
 \qquad
 \zeta(r)-H_n^{(r)}=
       \frac{(-1)^r}{(r-1)!}\psi^{(r-1)}(n+1).
\]
At every fixed truncation order $J$, the polynomial substitution has
remainder
\[
 O\bigl(x^{-J-1}(1+\log^d x)\bigr)
\]
for some fixed $d$.
Subtract those finitely many terms from the Dirichlet series. The
remaining series converges normally in $\Re s>-J$ on compact subsets.
The subtracted terms are
\[
 \sum_{j=0}^J\sum_\ell
   q_{P,j,\ell}(a)(-1)^\ell\partial_s^\ell\zeta(s+j,a).
\]
Their only principal contribution at $s=1-j$ is
\eqref{ghp:eq:principal}. Taking arbitrarily large $J$ proves the
continuation. Distinct powers $(s+j-1)^{-\ell-1}$ cannot cancel one
another, proving the last assertion.
\end{proof}

At $a=1/2$, all odd Bernoulli polynomials vanish. Consequently
\begin{equation}\label{ghp:eq:formal-parity}
 T_1(-z,L;1/2)=T_1(z,L;1/2),\qquad
 T_r(-z;1/2)=\begin{cases}
  2\zeta(r)-T_r(z;1/2),&r\text{ even},\\
  T_r(z;1/2),&r\text{ odd}.
 \end{cases}
\end{equation}
This proves the sufficiency of $P=\tau P$. The necessity requires
showing that no nonzero polynomial is annihilated by the formal
substitution. A finite list of verified asymptotic coefficients could
not establish that fact.

\subsection{Why the formal substitution is injective}

We use the following established result in exactly its formal Laurent
setting. Adamczewski, Dreyfus, and Hardouin
\cite[Theorem~1.2, case $\mathbf S_\infty$]{ghp:ADH} prove that a formal
Laurent series $f\in\mathbb C((x^{-1}))$ satisfying a homogeneous linear
difference equation over $\mathbb C(x)$ for the shift $x\mapsto x+1$
is either rational or differentially transcendental over
$\mathbb C(x)$. The latter means that every finite tuple
$f,f',\ldots,f^{(q)}$ is algebraically independent over $\mathbb C(x)$.

\begin{lemma}[Independence of the centered formal harmonic tails]
\label{ghp:lem:injective}
The substitution
\[
 \mathbb C[X_1,\ldots,X_R]\longrightarrow\mathbb C[L][[z]],
 \qquad X_r\longmapsto T_r(z,L;1/2),
\]
is injective.
\end{lemma}

\begin{proof}
Let
\begin{equation}\label{ghp:eq:formal-trigamma}
 f(x)=\frac1x+\sum_{k\geq1}B_{2k}(1/2)x^{-2k-1}
        \in\mathbb C((x^{-1})).
\end{equation}
This is the formal expansion of $\psi'(x+1/2)$. The trigamma shift
identity, or the Bernoulli generating series, gives
\begin{equation}\label{ghp:eq:shift}
 f(x+1)-f(x)=-\frac1{(x+1/2)^2}.
\end{equation}
In particular it satisfies the homogeneous second-order equation
\begin{equation}\label{ghp:eq:homogeneous}
 f(x+2)-(1+c(x))f(x+1)+c(x)f(x)=0,
 \qquad c(x)=\frac{(x+1/2)^2}{(x+3/2)^2}.
\end{equation}

There is no rational solution of \eqref{ghp:eq:shift}. Indeed, in any
nonempty finite set of poles of a rational function belonging to one
coset modulo $\mathbb Z$, choose the first and last pole in the integer
ordering. Its forward difference has an uncancelled pole immediately
before the first and another at the last. These poles are distinct.
Thus a rational forward difference cannot have precisely one pole in
such a coset. The right side of \eqref{ghp:eq:shift} has exactly one
pole. A rational function with no finite poles is a polynomial, whose
forward difference is polynomial, and is also impossible.

The cited difference-equation theorem applied to
\eqref{ghp:eq:homogeneous} now shows that
$f,f',\ldots,f^{(R-2)}$ are algebraically independent over
$\mathbb C(x)$ when $R\geq2$. For $r\geq2$, the substitution for $X_r$
is the affine expression
\[
 T_r=\zeta(r)-\frac{(-1)^r}{(r-1)!}f^{(r-2)}.
\]
These tails are therefore algebraically independent, even over
$\mathbb C(x)$. Finally $T_1=L+\gamma+A(z,1/2)$ has coefficient one
in the independent formal symbol $L$, and the other $T_r$ contain no
$L$. If a nonzero polynomial in $X_1$ has degree $d$ and nonzero leading
coefficient $a_d(X_2,\ldots,X_R)$, the coefficient of $L^d$ after
substitution is the nonzero substituted value of $a_d$. This proves
injectivity, including the case $R=1$ directly.
\end{proof}

\begin{proof}[Proof of Theorem~\ref{ghp:thm:classification}]
By Proposition~\ref{ghp:prop:principal}, regularity at every $s=-2m$
is equivalent to the vanishing of every odd coefficient $q_{P,2m+1}$.
Equivalently, the complete formal substitution is even in $z$ while
$L$ is held fixed. By \eqref{ghp:eq:formal-parity}, this says that the
substitution annihilates $P-\tau P$. Lemma~\ref{ghp:lem:injective}
therefore gives $P-\tau P=0$.

Writing the even-index variables as $\zeta(2j)-V_j$ makes $\tau$ the
simultaneous sign change of all the $V_j$. The invariant polynomials
are exactly those with even total degree in these variables. Every
such monomial can be grouped into pairs $V_iV_j$, proving the stated
finite generating description.
\end{proof}

\begin{corollary}[Complementary pole lattices]\label{ghp:cor:split}
Write $P=P_++P_-$, where $P_\pm=(P\pm\tau P)/2$. Then
$\mathcal D_{P_+}(s,1/2)$ has possible poles only at
$1,-1,-3,\ldots$, while $\mathcal D_{P_-}(s,1/2)$ has possible poles
only at $0,-2,-4,\ldots$. Moreover, if $P\ne0$,
$\mathcal D_P(s,1/2)$ is not entire.
\end{corollary}
\begin{proof}
The corresponding formal expansions have only even and only odd
powers of $z$, respectively. This gives the two pole lists by
\eqref{ghp:eq:principal}. If the whole Dirichlet function were entire,
all its coefficient polynomials $q_{P,j}$ would vanish. Injectivity
would then force $P=0$.
\end{proof}

This is a statement about identities of functions and formal series,
not about algebraic independence of their special numerical values.
The criterion concerns cancellation at \emph{every} point of the
specified parity. At any finite list of spectral points, additional
accidental cancellations can occur and are decided by the finite
Bernoulli conditions \eqref{ghp:eq:principal}.

\section{All centered trigamma-square moments and a polylogarithmic generator}
\label{ghp:square:sec}

The invariant numerator $(\zeta(2)-H_n^{(2)})^2$ provides an explicit
family beyond raw powers of $H_n$. Define
\begin{equation}\label{ghp:def:square}
 A_n=\zeta(2)-H_n^{(2)}=\psi'(n+1),\qquad
 \mathcal T(s)=\sum_{n=0}^{\infty}\frac{A_n^2}{(n+1/2)^s}.
\end{equation}
The defining series already converges for $\Re s>-1$. By the preceding
theorem it is regular at every nonpositive even integer. All its values
there reduce to a rational number plus a rational multiple of
$\zeta(3)$.

\begin{theorem}[Every nonpositive even trigamma-square value]
\label{ghp:thm:square}
For every integer $r\geq0$,
\begin{align}\label{ghp:eq:all-square}
 \mathcal T(-2r)={}&3B_{2r}(1/2)\zeta(3)\notag\\
 &-\frac12\sum_{j=1}^{r}\binom{2r}{2j}
   \frac{2j-1}{2j+1}B_{2r-2j}(1/2)B_{2j-2}.
\end{align}
The Bernoulli number convention is $B_1=-1/2$, and the sum is empty
when $r=0$.
\end{theorem}

For example,
\begin{align}\label{ghp:eq:square-table}
 \mathcal T(0)&=3\zeta(3),\notag\\
 \mathcal T(-2)&=-\frac14\zeta(3)-\frac16,\notag\\
 \mathcal T(-4)&=\frac7{80}\zeta(3)+\frac1{30},\notag\\
 \mathcal T(-6)&=-\frac{31}{448}\zeta(3)+\frac1{672},\notag\\
 \mathcal T(-8)&=\frac{127}{1280}\zeta(3)-\frac{17}{540}.
\end{align}
The first value is the classical Euler sum. The theorem supplies a
single finite formula for all subsequent centered moments.

\subsection{Equivalent identities of absolutely convergent ordinary sums}

Let
\begin{equation}\label{ghp:def:ck}
 c_k=\sum_{i=0}^{k-1}B_{2i}(1/2)B_{2k-2-2i}(1/2),\qquad k\geq1.
\end{equation}
Thus $c_1=1$, $c_2=-1/6$, $c_3=47/720$, and
$c_4=-257/5040$. Squaring \eqref{ghp:eq:formal-trigamma} gives
$A_n^2\sim\sum_{k\geq1}c_k(n+1/2)^{-2k}$.
For every $r\geq0$ the exact ordinary convergent identity is
\begin{equation}\label{ghp:eq:convergent-square}
 \sum_{n=0}^{\infty}\left\{(n+1/2)^{2r}\psi'(n+1)^2
       -\sum_{k=1}^r c_k(n+1/2)^{2r-2k}\right\}
       =\mathcal T(-2r),
\end{equation}
where the right side is explicitly \eqref{ghp:eq:all-square}. The
summands after subtraction are $O(n^{-2})$.

For instance,
\[
 \sum_{n=0}^{\infty}\left((n+1/2)^2\psi'(n+1)^2-1\right)
       =-\frac14\zeta(3)-\frac16,
\]
and
\[
 \sum_{n=0}^{\infty}\left((n+1/2)^4\psi'(n+1)^2
                         -(n+1/2)^2+\frac16\right)
       =\frac7{80}\zeta(3)+\frac1{30}.
\]
These statements need no prescription for summing a divergent series.
To prove \eqref{ghp:eq:convergent-square}, subtract the displayed
asymptotic terms in the half-plane of convergence and continue. Their
Dirichlet sums at $s=-2r$ are
$c_k\zeta(2k-2r,1/2)=0$, since
$\zeta(-2m,1/2)=0$ for $m\geq0$.

\subsection{Proof of the Bernoulli formula by finite summation}

For a sequence admitting an expansion in powers of $N$ and $\log N$,
write $\operatorname{CT}_N$ for its coefficient of $N^0(\log N)^0$.
Only finite exact identities followed by known asymptotic expansions
will be used here. Put
\begin{equation}\label{ghp:def:S}
 S_r(k)=\sum_{n=0}^{k-1}(n+1/2)^{2r}
  =\frac{B_{2r+1}(k+1/2)}{2r+1}
  =\sum_{j=0}^r b_{r,j}k^{2j+1},
\end{equation}
where
\begin{equation}\label{ghp:def:b}
 b_{r,j}=\frac{\binom{2r}{2j}}{2j+1}B_{2r-2j}(1/2).
\end{equation}
The Bernoulli polynomial $B_{2r+1}(1/2)$ is zero, so this is an odd
polynomial in $k$ and has zero constant term.

\begin{proof}[Proof of Theorem~\ref{ghp:thm:square}]
First, \eqref{ghp:eq:convergent-square} and \eqref{ghp:def:S} show that
\begin{equation}\label{ghp:eq:ct-reduction}
 \mathcal T(-2r)=\operatorname{CT}_N
       \sum_{n=0}^{N-1}(n+1/2)^{2r}A_n^2.
\end{equation}
Indeed every subtracted even power has partial sum with zero constant
term in $N$, and the convergent remainder tends to the stated value.

The identity $A_{k-1}=A_k+k^{-2}$ gives finite summation by parts:
\begin{align}\label{ghp:eq:partial-sum}
 \sum_{n=0}^{N-1}(n+1/2)^{2r}A_n^2
 ={}&\sum_{k=1}^N S_r(k)
          \left(\frac{2A_k}{k^2}+\frac1{k^4}\right)
       +S_r(N)A_N^2.
\end{align}
The coefficient $b_{r,0}$ of $k$ contributes
\[
 \operatorname{CT}_N
 \left(2\sum_{k=1}^N\frac{A_k}{k}+H_N^{(3)}+NA_N^2\right)
 =3\zeta(3).
\]
Here $\sum_{k\geq1}A_k/k=\zeta(2,1)=\zeta(3)$ is the classical
Euler identity and $NA_N^2\to0$. For completeness, the Euler identity
follows from
\[
 \zeta(2,1)=\frac12\int_0^1\frac{\log^2(1-t)}t\,dt
 =\frac12\int_0^\infty\frac{u^2}{e^u-1}\,du=\zeta(3).
\]

For $j\geq1$, it remains to compute the constant of
\begin{equation}\label{ghp:def:Ej}
 E_j(N)=2\sum_{k=1}^Nk^{2j-1}A_k
        +\sum_{k=1}^Nk^{2j-3}+N^{2j+1}A_N^2.
\end{equation}
Let
\[
 F_j(N)=\sum_{k=1}^Nk^{2j-1}
       =\frac{B_{2j}(N+1)-B_{2j}}{2j}.
\]
Changing the order in a finite sum gives
\begin{equation}\label{ghp:eq:finite-change-order}
 \sum_{k=1}^Nk^{2j-1}A_k
  =A_NF_j(N)+\sum_{\ell=1}^N\frac{F_j(\ell-1)}{\ell^2}.
\end{equation}
For $j\geq2$, the second summand is a sum of polynomial powers of
$\ell$ and has zero constant term as a polynomial in $N$;
$B_{2j-1}=0$ removes a possible $1/\ell$ term. The other power sum in
\eqref{ghp:def:Ej} also has zero constant. For $j=1$ the second sum in
\eqref{ghp:eq:finite-change-order} equals $(N-H_N)/2$; its harmonic
term cancels exactly against $\sum k^{-1}$ in \eqref{ghp:def:Ej}.
In all cases,
\begin{equation}\label{ghp:eq:Ej-CT}
 \operatorname{CT}_NE_j(N)
   =2\operatorname{CT}_N(A_NF_j(N))
          +\operatorname{CT}_N(N^{2j+1}A_N^2).
\end{equation}

Now use the uncentered tail expansion
\[
 A_N=\frac1N-\frac1{2N^2}
                  +\sum_{k\geq1}\frac{B_{2k}}{N^{2k+1}}.
\]
The finite coefficient calculation gives
\begin{equation}\label{ghp:eq:two-constants}
 \operatorname{CT}_N(A_NF_j(N))
      =\frac{3-2j}{4}B_{2j-2},\qquad
 \operatorname{CT}_N(N^{2j+1}A_N^2)=-B_{2j-2}.
\end{equation}
To check the first identity, the coefficient of $N^\ell$ in $F_j$
is $\binom{2j}{\ell}B_{2j-\ell}(1)/(2j)$. Only $\ell=2$ and
$\ell=2j-1$ can pair with nonzero inverse coefficients of $A_N$
when $j\geq2$; their contributions are
$-(2j-1)B_{2j-2}/4$ and $B_{2j-2}/2$. The case $j=1$ is the direct
product with $F_1(N)=(N^2+N)/2$ and gives $1/4$.
For the second identity, an odd inverse exponent in $A_N^2$ arises
only by multiplying $-1/(2N^2)$ by the appropriate odd inverse power;
this gives $-B_{2j-2}$, also at $j=1$.

Combining \eqref{ghp:eq:Ej-CT} and \eqref{ghp:eq:two-constants} yields
\[
 \operatorname{CT}_NE_j(N)=-\frac{2j-1}{2}B_{2j-2}.
\]
Multiplication by $b_{r,j}$, followed by the $j=0$ contribution,
proves \eqref{ghp:eq:all-square}.
\end{proof}

\subsection{A convergent generator involving classical polylogarithms}

The complete sequence has an analytic exponential generating function.
Let
\[
 \mathcal G(t)=\sum_{r=0}^{\infty}
               \mathcal T(-2r)\frac{t^{2r}}{(2r)!}.
\]

\begin{theorem}[Polylogarithmic moment generator]\label{ghp:thm:generator}
The series $\mathcal G$ converges for $|t|<2\pi$ and equals
\begin{equation}\label{ghp:eq:integral-generator}
 \mathcal G(t)=\frac{t}{2\sinh(t/2)}
   \left[3\zeta(3)-\frac1{4t}
       \int_0^t(t-u)u^2\coth(u/2)\,du\right].
\end{equation}
The singularities at zero in this expression are removable. For
real $0<t<2\pi$, an equivalent formula in classical polylogarithms is
\begin{align}\label{ghp:eq:polylog-generator}
 \mathcal G(t)=\frac{t}{2\sinh(t/2)}\bigg[
 2\zeta(3)-\frac{t^3}{48}
 -\frac t2\operatorname{Li}_2(e^{-t})
 -2\operatorname{Li}_3(e^{-t})
 +\frac{3\bigl(\zeta(4)-\operatorname{Li}_4(e^{-t})\bigr)}t
 \bigg].
\end{align}
Its continuation specified by \eqref{ghp:eq:integral-generator} is
holomorphic and even throughout $|t|<2\pi$.
\end{theorem}

\begin{proof}
The Bernoulli generating functions give
\[
 \frac{t}{2\sinh(t/2)}
     =\sum_{r\geq0}B_{2r}(1/2)\frac{t^{2r}}{(2r)!},
 \qquad
 \frac u2\coth(u/2)
     =\sum_{k\geq0}B_{2k}\frac{u^{2k}}{(2k)!}.
\]
Both converge in the disk of radius $2\pi$. Integrating the second
series twice gives
\[
 \frac1{2t}\int_0^t(t-u)u^2\coth(u/2)\,du
  =\sum_{j\geq1}\frac{2j-1}{2j+1}B_{2j-2}
                  \frac{t^{2j}}{(2j)!}.
\]
Thus \eqref{ghp:eq:integral-generator} is exactly the binomial
convolution in \eqref{ghp:eq:all-square}. This proves convergence,
analyticity and the integral formula. In particular all apparent odd
or logarithmic terms in any branch representation must cancel.

For positive $t$, expand $(e^u-1)^{-1}$ in its convergent positive
exponential series. Direct integration, justified for the integrable
factor $(t-u)u^2$, gives
\begin{align*}
 \int_0^t\frac{(t-u)u^2}{e^u-1}\,du
  ={}&2t\zeta(3)-6\zeta(4)
      +t^2\operatorname{Li}_2(e^{-t})\\
     &+4t\operatorname{Li}_3(e^{-t})
      +6\operatorname{Li}_4(e^{-t}).
\end{align*}
Finally $\coth(u/2)=1+2/(e^u-1)$ and
$\int_0^t(t-u)u^2\,du=t^4/12$ give
\eqref{ghp:eq:polylog-generator}.
\end{proof}

\subsection{Verification, audit scope, and further exact questions}

The accompanying \texttt{verify\_harmonic\_parity.py} implements the
formal involution, verifies the centered trigamma difference equation,
and checks the two finite Bernoulli constant calculations independently
for eight values of $j$. It evaluates seven spectral values directly
from finite trigamma sums with independently truncated asymptotic tails
at two cutoffs, using 110 decimal digits. The largest discrepancy
against \eqref{ghp:eq:all-square} is less than $1.6\times10^{-85}$.
At three positive parameters it compares the polylogarithmic generator
both with the Bernoulli moment series and with ordinary quadrature of
\eqref{ghp:eq:integral-generator}. These are floating-point diagnostics,
not interval enclosures. The proofs do not depend on their precision.

The selected harmonic statements in the pinned report were consistent
with these checks. No new false formula was identified in that selected
source. Its Q11 warning was correct: raw-power parity alone does not
apply to an arbitrary higher-harmonic numerator. The result here
resolves that specific question by supplying the exact invariant ring;
it does not label the earlier warning an error.

The following are concrete extensions beyond the proved results.
\begin{enumerate}
\item Evaluate every centered moment of
$(\zeta(2i)-H_n^{(2i)})(\zeta(2j)-H_n^{(2j)})$ by a uniform finite
Bernoulli formula, and determine the precise convergent multiple-zeta
coordinates that remain. The current trigamma-square result is
$i=j=1$.
\item For invariant products of three or more tail pairs, seek a
polylogarithmic generator analogous to \eqref{ghp:eq:polylog-generator}
and determine when the resulting multiple zeta values reduce to
ordinary zeta products.
\item Classify cancellation at a \emph{finite} prescribed set of
spectral integers by the finite coefficient map
\eqref{ghp:eq:principal}. Its kernel can be larger than the invariant
ring and is a tractable exact linear problem in any bounded polynomial
space.
\item Derive exact spectral-derivative identities for $\mathcal T$
at its regular even points. Such derivatives retain a global remainder
and do not follow by formally differentiating the discrete index $r$
in \eqref{ghp:eq:all-square}.
\item Determine an elementary proof of the injectivity lemma tailored
to the additive trigamma difference equation. This would remove the
use of the general difference Galois theorem, while leaving the current
unconditional classification unchanged.
\end{enumerate}

% END sections/02_harmonic_parity.tex


% BEGIN sections/03_resolvent.tex
\section{Cotangent resolvents: a nonpolynomial Stieltjes generator}
\label{res:section}

Question Q8 of the incoming \emph{Exact Resonant Identities} report
\cite{ghp:ExactResonant} asks for a nonpolynomial extension of its
cotangent moment calculus, including the local subtraction before
parameter expansion. We give such an extension for every rational
function of the cotangent bounded at infinity and with no real pole. A single
resolvent supplies the answer; differentiation gives all repeated
poles, and partial fractions give the general rational case. The
answer involves ordinary polylogarithms and their order derivatives.
The Fourier identity used in the proof is classical. The contribution
here is the completed rational-kernel identity, its precise endpoint
correction at every Stieltjes and argument-derivative order, and its
compatibility with continuous polygamma orders and primitives.

\subsection{Coordinates, branches, and finite parts}

Write
\begin{equation}\label{res:kernel}
 K(x)=\pi\cot(\pi x),\qquad
 R_z(x)=\frac1{K(x)-z},\qquad 0<x<1,\quad z\in\mathbb C\setminus\mathbb R.
\end{equation}
Although $K$ has endpoint poles, $R_z$ extends to a smooth periodic
function, with $R_z(0)=R_z(1)=0$. On either open half-plane put
\begin{equation}\label{res:coordinates}
 \epsilon=\operatorname{sgn}(\Im z),\quad
 a=z+\epsilon\pi i,\quad
 q=\frac{z-\epsilon\pi i}{z+\epsilon\pi i},\quad
 c=\frac{2\epsilon\pi i}{a^2},\quad
 \sigma=-2\epsilon\pi i.
\end{equation}
Thus $|q|<1$, and throughout the section
\[
 \Log\sigma=\log(2\pi)-\epsilon i\pi/2.
\]
The letter $a$ in this section is only the denominator in
\eqref{res:coordinates}; it is not a Hurwitz endpoint.
To remove an otherwise unnecessary exceptional case, define
\begin{equation}\label{res:normalizedpolylog}
 \Phi_*(s,q)=\frac{\Li_s(q)}q
       =\sum_{n\ge1}\frac{q^{n-1}}{n^s},\qquad
 \Phi_*(s,0)=1.
\end{equation}
This function is entire in $s$ and holomorphic for $|q|<1$.
Its order derivatives at $q=0$ are zero except at order zero.

We use the conventional generalized Stieltjes functions
\begin{equation}\label{res:stieltjes}
 \zeta(1-s,x)=-\frac1s+
       \sum_{m\ge0}\gamma_m(x)\frac{s^m}{m!},\qquad
 \gamma_0(x)=-\psi(x).
\end{equation}
A superscript $(p)$ on $\gamma_m$ or $\psi$ means an argument
derivative. A superscript $[j]$ on $\Li_s$ or $\Phi_*$ means an order
derivative. These two operations must be distinguished.

Every finite part below is taken in the original variable $x$:
$\operatorname{FP}_x\int_0^1 f(x)\,dx$ is the coefficient of
$\varepsilon^0(\log\varepsilon)^0$ in
$\int_\varepsilon^{1-\varepsilon}f(x)\,dx$. Equivalently, subtract
the finitely many nonintegrable endpoint terms, integrate the
remainder, and add the constant terms of the explicit subtractions.
For the resolvent integrals only the endpoint $x=0$ can be singular.
The Stieltjes shift formula gives
\begin{equation}\label{res:singularpart}
 \gamma_m(x)=\frac{\log^m x}{x}+\gamma_m(x+1).
\end{equation}
Consequently all the finite parts used here exist. In particular,
\[
 M_{m,p}(z):=\operatorname{FP}_x\int_0^1\gamma_m^{(p)}(x)R_z(x)\,dx
\]
is an ordinary convergent integral when $p=0$.

\subsection{The complete all-index identity}

Set
\begin{equation}\label{res:BE}
 B_\epsilon(s)=\Gamma(1+s)\sigma^{-s},\qquad
 E_p(s)=\prod_{j=1}^p\left(1-\frac{s}{j}\right),\qquad E_0(s)=1.
\end{equation}

\begin{theorem}[Completed cotangent--Stieltjes resolvent]
\label{res:main}
For all integers $m,p\ge0$ and $z\notin\mathbb R$,
\begin{equation}\label{res:master}
 \boxed{\displaystyle
 M_{m,p}(z)=m!c\sigma^p[s^{m+1}]
 \left\{B_\epsilon(s)\Phi_*(s-p,q)
                   -E_p(s)\Phi_*(-p,q)\right\}.}
\end{equation}
The expression in braces vanishes at $s=0$. The identity is
holomorphic in $z$ on each half-plane, including $z=\epsilon i\pi$.
For $p\ge1$ the polynomial $E_p$ is the full local endpoint
correction. Its contribution vanishes whenever $m\ge p$.
\end{theorem}

The identity can be expanded without any coefficient notation.
Let $g_j=B_\epsilon^{(j)}(0)$ and $e_{p,j}=E_p^{(j)}(0)$. Then
\begin{equation}\label{res:expanded}
 M_{m,p}(z)=\frac{c\sigma^p}{m+1}
 \left\{\sum_{j=0}^{m+1}\binom{m+1}{j}
     g_{m+1-j}\Phi_*^{[j]}(-p,q)
                 -e_{p,m+1}\Phi_*(-p,q)\right\}.
\end{equation}
The coefficients $g_j$ are complete exponential Bell polynomials
with logarithmic derivatives
\[
 b_1=-\gamma-\Log\sigma,\qquad
 b_j=(-1)^j(j-1)!\zeta(j)\quad(j\ge2),
\]
so $g_0=1$, $g_1=b_1$, $g_2=b_1^2+\zeta(2)$, and
$g_3=b_1^3+3b_1\zeta(2)-2\zeta(3)$.
Likewise
\begin{equation}\label{res:contactcoeff}
 e_{p,j}=(-1)^j j!\,
 e_j\left(1,\frac12,\ldots,\frac1p\right),
\end{equation}
where $e_j$ on the right is the elementary symmetric polynomial;
it is zero for $j>p$. Thus the contact terms use ordinary harmonic
numbers and their finite Bell-polynomial combinations.

\begin{proof}[Proof of Theorem~\ref{res:main}]
It is enough to prove the upper-half-plane case; reversing the
Fourier signs gives the other case with exactly the branches in
\eqref{res:coordinates}. Put $w=e^{2\pi ix}$. Direct algebra gives
the uniformly convergent Fourier expansion
\begin{equation}\label{res:fourierkernel}
 R_z(x)=-\frac1a+c\sum_{n\ge1}q^{n-1}e^{2\pi inx}.
\end{equation}
For example, at $z=i\pi$ this reduces to
$(e^{2\pi ix}-1)/(2\pi i)$. The classical Hurwitz Fourier formula
\cite[\S25.11]{DLMF} implies, initially in a sufficiently large
right half-plane,
\[
 \int_0^1\zeta(1-s,x)\,dx=0,\qquad
 \int_0^1\zeta(1-s,x)e^{2\pi inx}\,dx
       =\Gamma(s)\sigma^{-s}n^{-s}.
\]
All termwise integrations are legitimate there because of the
geometric factor $q^n$. Argument differentiation multiplies the
negative Fourier mode by $\sigma n$. For $\Re s>p+1$, this gives
\begin{equation}\label{res:spectralintegral}
 J_p(s,z):=\int_0^1\partial_x^p\zeta(1-s,x)R_z(x)\,dx
       =c\Gamma(s)\sigma^{p-s}\Phi_*(s-p,q).
\end{equation}
This equation is then an identity of meromorphic continuations.

We now relate that continuation to the pointwise finite part;
this is the step that cannot be replaced by substitution at $s=0$.
Write $r_j(z)=[x^j]R_z(x)$. For $p\ge1$, the sole local term
whose moving exponent crosses $-1$ at $s=0$ is
\begin{equation}\label{res:localmoving}
 r_p(z)A_p(s)x^{s-1},\qquad
 A_p(s)=\prod_{j=1}^p(s-j).
\end{equation}
Indeed the singular part of
$\partial_x^p\zeta(1-s,x)$ is
$A_p(s)x^{s-p-1}$. The analytic part from $\zeta(1-s,x+1)$
has no pole in $s$ after an argument derivative. Every other
Taylor term of $R_z$ has an exponent separated from $-1$ near
zero. For such a term, finite-part integration commutes with
coefficient extraction. The term \eqref{res:localmoving} is different:
its analytically continued integral is $r_pA_p(s)/s$, whereas the
finite part of every coefficient $x^{-1}\log^j x$ on $(0,1)$ is zero.
It follows, with the entire amplitude retained, that
\begin{equation}\label{res:completion}
 M_{m,p}(z)=m![s^m]
            \left(J_p(s,z)-\frac{r_p(z)A_p(s)}s\right),\quad p\ge1.
\end{equation}
One can make this argument on $(0,\delta)$ with a finite Taylor
subtraction and a normally integrable remainder. The compensating
integrals on $(\delta,1)$ cancel the artificial splitting point.
This also proves all holomorphy and coefficient interchanges used here.

Differentiating \eqref{res:fourierkernel} at $x=0$ gives
\begin{equation}\label{res:rj}
 r_p(z)=\frac{c(-\sigma)^p}{p!}\Phi_*(-p,q),\qquad p\ge1.
\end{equation}
Since $A_p(s)=(-1)^pp!E_p(s)$, substitution in
\eqref{res:completion} proves \eqref{res:master} for $p\ge1$.

For $p=0$, the product $\gamma_m R_z$ is integrable by
\eqref{res:singularpart}. The pole in \eqref{res:stieltjes} is
independent of $x$, so
\[
 \sum_{m\ge0}M_{m,0}(z)\frac{s^m}{m!}
       =J_0(s,z)+\frac1s\int_0^1R_z(x)\,dx
       =J_0(s,z)-\frac1{as}.
\]
But $c\Phi_*(0,q)=1/a$. This is exactly
\eqref{res:master} with $E_0=1$. Thus the same formula accounts
for both the global Hurwitz pole at $p=0$ and the local moving
endpoint singularity at $p\ge1$.
\end{proof}

\subsection{Explicit examples and a check of the correction}

At the first Stieltjes index the formula is
\begin{equation}\label{res:polygamma}
 \operatorname{FP}_x\int_0^1\psi^{(p)}(x)R_z(x)\,dx
  =c\sigma^p\left\{(\gamma+\Log\sigma-H_p)\Phi_*(-p,q)
                                      -\Phi_*^{[1]}(-p,q)\right\},
\end{equation}
with $H_0=0$. The harmonic number is essential whenever $p>0$.
For $m=1$ the complete formula is
\begin{align}\label{res:gamma1}
 M_{1,p}(z)=\frac{c\sigma^p}{2}\bigl\{
  &\Phi_*^{[2]}(-p,q)-2(\gamma+\Log\sigma)\Phi_*^{[1]}(-p,q)\notag\\
  &+\bigl((\gamma+\Log\sigma)^2+\zeta(2)
                      -H_p^2+H_p^{(2)}\bigr)\Phi_*(-p,q)\bigr\}.
\end{align}
Here $H_0^{(2)}=0$. The correction in this equation vanishes at
$p=0,1$, as the general threshold predicts.

At $z=i\pi$, the normalized polylogarithm is identically one,
and all positive order derivatives disappear. In particular,
\begin{equation}\label{res:explicitcheck}
 \operatorname{FP}_x\int_0^1
       \frac{\psi'(x)}{\pi\cot(\pi x)-i\pi}\,dx
       =1-\gamma-\log(2\pi)+\frac{i\pi}{2}.
\end{equation}
There is a direct independent verification of its constant term.
At this $z$, $R'_z=e^{2\pi ix}$ and the finite boundary value of
$[\psi R_z]_0^1$ is $1$, because $\psi(x)R_z(x)\to-1$ as
$x\downarrow0$. Integration by parts therefore supplies precisely
the $1$ in \eqref{res:explicitcheck}. Omitting the local correction
would miss it.

The local coefficients themselves require no special functions:
\begin{equation}\label{res:riccati}
 R_z'=1+2zR_z+(z^2+\pi^2)R_z^2,\qquad R_z(0)=0.
\end{equation}
Thus $r_1=1$, $r_2=z$, $r_3=z^2+\pi^2/3$, and
$r_4=z^3+2\pi^2z/3$. In general $r_p$ is a polynomial in $z$
of degree $p-1$. The difference between the raw spectral coefficient
and the stated finite part is therefore a finite polynomial in $z$,
not an additional unevaluated integral.

For a real parameter $b>0$, taking the real and imaginary parts of
the $p=m=0$ formula at $z=ib$ gives ordinary real integrals. For example,
with $q=(b-\pi)/(b+\pi)$,
\begin{align}
 \int_0^1\frac{\gamma_0(x)K(x)}{K(x)^2+b^2}\,dx
       &=\frac{\pi}{2(b+\pi)},\label{res:realexample1}\\
 \int_0^1\frac{\gamma_0(x)}{K(x)^2+b^2}\,dx
       &=\frac{\gamma+\log(2\pi)
                   -(1-q)\Phi_*^{[1]}(0,q)}{b(b+\pi)}.
       \label{res:realexample2}
\end{align}
At $b=3\pi$, the last numerator becomes
$\gamma+\log(2\pi)-\Li_0^{[1]}(1/2)$.
Equation \eqref{res:realexample1} also follows directly from the
digamma reflection formula, providing a separate check of the phase.

\subsection{Repeated poles and the complete bounded-rational family}

\begin{corollary}[All repeated resolvents]\label{res:repeated}
For every integer $r\ge1$,
\begin{equation}\label{res:repeatedformula}
 \operatorname{FP}_x\int_0^1
       \frac{\gamma_m^{(p)}(x)}{(K(x)-z)^r}\,dx
       =\frac1{(r-1)!}\partial_z^{r-1}M_{m,p}(z).
\end{equation}
This is an ordinary absolutely convergent integral if $r>p$.
Every rational function $Q(u)$ bounded at infinity and with no real pole admits
a finite reduction of
$\operatorname{FP}_x\int_0^1\gamma_m^{(p)}(x)Q(K(x))\,dx$
to the right-hand sides of \eqref{res:repeatedformula}.
\end{corollary}

\begin{proof}
Parameter differentiation commutes with the finitely many local
subtractions on compact subsets away from the real poles. Since
$\partial_z^{r-1}R_z=(r-1)!R_z^r$, the formula follows.
Near zero the integrand has size
$O(x^{r-p-1}(1+|\log x|^m))$, proving the convergence assertion.
Equivalently, the resonant coefficient $[x^p]R_z^r$ is zero
when $r>p$, so all contact terms disappear from the differentiated
formula.

Partial fractions express $Q$ as a constant plus finitely many
repeated resolvents. The constant contributes zero:
\[
 \operatorname{FP}_x\int_0^1\gamma_m^{(p)}(x)\,dx=0.
\]
For $p=0$ this follows by integrating the Hurwitz spectral family
and matching its constant and endpoint pole. For $p\ge1$ it
follows from the finite part of the endpoint difference of
$\gamma_m^{(p-1)}$, using \eqref{res:singularpart}. The analytic
parts at $0$ and $1$ agree, and the singular part has zero constant.
This proves the complete finite reduction.
\end{proof}

Repeated $z$ derivatives introduce no new transcendental function
class: $\partial_q\Li_s(q)=\Li_{s-1}(q)/q$, and $q(z)$ is rational.
The reduction therefore stays within finitely many polylogarithmic
order derivatives at explicit arguments. This proves the rational
sector of the proposed nonpolynomial extension. It does not assert
the same finite closure for arbitrary complex powers of $K-z$.

\section{Continuous polygamma orders, Gamma primitives, and boundary values}
\label{res:primitive-section}

\subsection{An entire transform and its integer-order anomaly}

We use the established generalized polygamma function of Espinosa
and Moll \cite{EM}, denoted here by $\Psi(v,x)$ to distinguish it
from the ordinary digamma:
\begin{equation}\label{res:EM}
 \Psi(v,x)=e^{-\gamma v}\partial_v
       \left(e^{\gamma v}\frac{\zeta(1+v,x)}{\Gamma(-v)}\right).
\end{equation}
The removable values in this definition are understood. It is entire
in $v$, satisfies $\partial_x\Psi(v,x)=\Psi(v+1,x)$, and equals
$\psi^{(p)}(x)$ when $v=p\ge0$. These are properties of the
classical generalized function, not new definitions of polygamma.

\Needspace{18\baselineskip}
\begin{theorem}[Continuous-order resolvent and finite-part conversion]
\label{res:continuous}
For $\Re v<1$ and $z\notin\mathbb R$, the following ordinary
absolutely convergent integral holds:
\begin{equation}\label{res:continuousformula}
 \mathcal C(v,z):=\int_0^1\Psi(v,x)R_z(x)\,dx
   =c\sigma^v\left\{(\gamma+\Log\sigma)\Phi_*(-v,q)
                                         -\Phi_*^{[1]}(-v,q)\right\}.
\end{equation}
The right side is entire in $v$. At a nonnegative integer $p$,
its continuation and the pointwise finite part are related by
\begin{equation}\label{res:anomaly}
 \operatorname{FP}_x\int_0^1\psi^{(p)}(x)R_z(x)\,dx
    =\mathcal C(p,z)-(-1)^pp!H_p r_p(z),
\end{equation}
where the correction is zero for $p=0$ without defining $r_0$.
\end{theorem}

\begin{proof}
The singular part of \eqref{res:EM} is a linear combination of
$x^{-v-1}$ and $x^{-v-1}\log x$ with entire coefficient functions
of $v$. Multiplication by $R_z=O(x)$ gives locally uniform
integrability when $\Re v<1$. Thus the integral is holomorphic
in that half-plane, including the removable point $v=0$.

For $\Re v<0$, divide \eqref{res:spectralintegral} with $p=0$
by $\Gamma(s)$ and set $s=-v$. Applying
$e^{-\gamma v}\partial_v e^{\gamma v}$ yields
\eqref{res:continuousformula}. Holomorphic continuation gives
the entire stated ordinary-integral strip. Finally, compare the
right side at $v=p$ with \eqref{res:polygamma}, and use
$c\sigma^p\Phi_*(-p,q)=(-1)^pp!r_p(z)$.
This proves the exact conversion formula.
\end{proof}

The last statement is a concrete instance of noncommutation between
specializing a parameter and taking an endpoint finite part. It is
not a defect of the generalized polygamma function. At $p=1$ the
correction in \eqref{res:anomaly} is $+1$ for every nonreal $z$.
At $p=2$ it is $-3z$. These exact corrections are useful checks
when a spectral formula is integrated into a symbolic system.

Integration by parts gives a compatible differential identity for
the entire transform:
\begin{equation}\label{res:continuous-recurrence}
 \mathcal C(v+1,z)
      =-2z\mathcal C(v,z)-(z^2+\pi^2)\partial_z\mathcal C(v,z).
\end{equation}
Initially take $\Re v<0$, so the boundary term vanishes and
$\int_0^1\Psi(v,x)\,dx=0$, and use \eqref{res:riccati}.
Both sides are entire in $v$, which proves the identity everywhere.
For the pointwise finite-part moments
$F_p(z)=\operatorname{FP}_x\int\psi^{(p)}R_z$, the corresponding
identity instead reads
\begin{equation}\label{res:pointwise-recurrence}
 F_{p+1}(z)=-2zF_p(z)-(z^2+\pi^2)F_p'(z)
                         +(-1)^pp!r_{p+1}(z).
\end{equation}
The last term is exactly the finite part of $[\psi^{(p)}R_z]_0^1$.
Thus the derivative hierarchy, the endpoint subtraction, and the
polylogarithmic formula agree at every order.

\subsection{Every balanced Gamma primitive}

For $r\ge0$, let
\begin{equation}\label{res:balanced}
 \mathcal B_r(x)=\frac{\zeta'(-r,x)+H_r\zeta(-r,x)}{r!}.
\end{equation}
These are the standard balanced negapolygamma functions
$\Psi(-r-1,x)$ \cite{EM}. Direct differentiation gives
\[
 \mathcal B_0(x)=\log\Gamma(x)-\tfrac12\log(2\pi),\qquad
 \mathcal B_r'(x)=\mathcal B_{r-1}(x)\quad(r\ge1).
\]
The additive constants have therefore been specified completely.

\begin{corollary}[The rational-kernel primitive ladder]
\label{res:balancedcor}
For every $r\ge0$,
\begin{equation}\label{res:balancedmoment}
 \int_0^1\mathcal B_r(x)R_z(x)\,dx
    =c\sigma^{-r-1}
       \left\{(\gamma+\Log\sigma)\Phi_*(r+1,q)
                                      -\Phi_*^{[1]}(r+1,q)\right\}.
\end{equation}
All these integrals are ordinary convergent integrals.
\end{corollary}

This follows from \eqref{res:continuousformula} at $v=-r-1$.
There is also a useful direct check. Differentiate
\eqref{res:spectralintegral} in $s$ at $s=r+1$. The derivative of
$\Gamma(s)$ supplies $H_r-\gamma$; the $H_r\zeta(-r,x)$ term
in \eqref{res:balanced} cancels $H_r$ exactly. That cancellation
explains why the same coefficient $\gamma+\Log\sigma$ works at
every primitive level.

In particular,
\begin{equation}\label{res:loggammamoment}
 \int_0^1\log\Gamma(x)R_z(x)\,dx
 =-\frac{\log(2\pi)}{2a}
   +\frac{\Phi_*^{[1]}(1,q)
                 -(\gamma+\Log\sigma)\Phi_*(1,q)}{a^2}.
\end{equation}
At the next level, the standard Barnes identity
\cite{Vardi}
\[
 \zeta'(-1,x)=\zeta'(-1)-\log G(x)+(x-1)\log\Gamma(x)
\]
shows that \eqref{res:balancedmoment} is an explicit formula for
the corresponding normalized Barnes-$G$ combination, since
$\zeta(-1,x)=-B_2(x)/2$. Here $G$ is the standard Barnes function, with its usual asymptotic
normalization, $G(x+1)=\Gamma(x)G(x)$, and $G(1)=1$. The
functional equation and the value at $1$ alone would not exclude
a multiplicative periodic factor. This supplies an exact primitive
family relevant to Q9 in the source, without claiming to evaluate
every unbalanced product containing $G$.

\subsection{Boundary values and recovery of the Stieltjes functions}

The resolvent also permits a useful inversion check. Given
$u\in\mathbb R$, let $x_u\in(0,1)$ be the unique point with
$K(x_u)=u$. For $m\ge0$ the ordinary moment $M_{m,0}$ has
one-sided boundary values satisfying
\begin{equation}\label{res:jump}
 M_{m,0}(u+i0)-M_{m,0}(u-i0)
       =\frac{2\pi i}{u^2+\pi^2}\gamma_m(x_u).
\end{equation}
Indeed the substitution $y=K(x)$ gives the Cauchy integral
\[
 M_{m,0}(z)=\int_{-\infty}^{\infty}
      \frac{\gamma_m(x_y)}{(y^2+\pi^2)(y-z)}\,dy.
\]
Its tail is integrable after the factor $(y-z)^{-1}$ is included,
and its density is smooth locally at $y=u$. Splitting off that
local value and integrating $1/(y-z)$ proves
\eqref{res:jump}; this is the elementary Cauchy jump formula.
The average of the two boundary values gives the corresponding
principal-value integral at the interior pole.

On the boundary, $q=e^{-2\pi ix_u}$ in the upper half-plane and
$q=e^{2\pi ix_u}$ in the lower one. Thus the two versions of
\eqref{res:master} recover the classical Hurwitz--polylogarithm
connection and all its Stieltjes order derivatives. This boundary
interpretation is a consistency check and an exact inversion formula;
it is not claimed as a new Hurwitz functional equation.

\subsection{Verification and the precise new scope}

The program \src{code/verify_resolvent.py} checks the Fourier and
Riccati expressions for the local coefficients by exact rational
function arithmetic, including both the Gamma polynomial and its
harmonic correction. Its numerical checks compare the formulas with
direct quadrature of digamma and generalized Stieltjes functions,
local-subtraction quadrature for the finite parts, convergent
integrals involving $\gamma_1'$, and the first three balanced
Gamma primitives. Complex-order tests use the defining Hurwitz
formula \eqref{res:EM}. The numerical and exact records are distinct.

The all-index rational closure answers the specified rational sector
of the source's nonpolynomial question. The local correction
mechanism extends the repository's established finite-part calculus;
it is not presented as an independently invented regularization.
The continuous-order conversion and the primitive ladder explain
which explicit boundary terms must accompany the resulting identities.

% END sections/03_resolvent.tex


% BEGIN sections/04_harmonic_reduction.tex
% Manuscript-ready module. Requires amsmath, amssymb, theorem environments.
% Source status: deductions beyond the two pinned incoming reports;
% classical Faulhaber, Hurwitz, Lerch and polynomial integration inputs
% are explicitly identified. No literature-wide novelty claim is made.

\section{Nonpositive harmonic indices and a convergent Lerch generator}
\label{harm:section}

The incoming independent-order report reduces a depth-two word whose
last index is a nonpositive integer. We extend that calculation to an
arbitrary position and depth, prove a canonical finite reduction, and
sum all zero-index decorations around one free index in a convergent
Lerch generating function. The summation of powers by Bernoulli
polynomials is classical; the contribution here is its compatibility
with the entire relative harmonic interpolant at arbitrary endpoints,
including its independent spectral derivatives. The interpolant itself
is established in the incoming report, and its one-endpoint version is
the interpolant of Ihara, Nakamura, and Yamamoto
\cite{harm:INY}. We do not claim literature-wide priority for these
consequences.

\subsection{Normalization and a local elimination theorem}

For $\Re a,\Re b>0$, use the normalization
\begin{align}
 Z_a(s_1,\ldots,s_r)
  &=\sum_{n_1>\cdots>n_r\ge0}
                  \prod_{j=1}^r(n_j+a)^{-s_j},\notag\\
 \mathcal H_{a,b}(w)
  &=Z_b(w)-\sum_{j=1}^{|w|}Z_a(w_{\le j})
                              \mathcal H_{a,b}(w_{>j}),
 \qquad \mathcal H_{a,b}(\varnothing)=1.
 \label{harm:normalization}
\end{align}
The $Z_a$ are meromorphically continued from their convergence domain.
The combined function $\mathcal H_{a,b}$ is entire in all spectral
indices. For $a=b+N$, $N\ge0$ an integer, it is the strict finite sum
over $N>n_1>\cdots>n_r\ge0$. In particular,
\begin{equation}
 h_{a,b}(s):=\mathcal H_{a,b}(s)
      =\zeta(s,b)-\zeta(s,a),\qquad
 h_{a,b}(1)=\psi(a)-\psi(b).
 \label{harm:h}
\end{equation}
We take $B_n=B_n(0)$, so $B_1=-1/2$, whereas $B_1(1)=1/2$.
This distinction is essential for strict inequalities.

\begin{theorem}[Eliminating a fixed nonpositive index]
\label{harm:local}
Let $m\ge0$ and $q=m+1$. All other indices below are arbitrary
independent complex numbers. With unchanged prefix and suffix words
suppressed, an interior index obeys
\begin{align}
 \mathcal H_{a,b}(\ldots,s,-m,t,\ldots)
  =\frac1q\sum_{k=0}^q\binom qk\bigl\{
   &B_{q-k}\mathcal H_{a,b}(\ldots,s-k,t,\ldots)\notag\\[-2pt]
   &-B_{q-k}(1)\mathcal H_{a,b}(\ldots,s,t-k,\ldots)
   \bigr\}.
 \label{harm:interior}
\end{align}
At the two boundaries the formulas are
\begin{align}
 \mathcal H_{a,b}(-m,t,\ldots)
 &=\frac1q\left\{B_q(a)\mathcal H_{a,b}(t,\ldots)
  -\sum_{k=0}^q\binom qk B_{q-k}(1)
                        \mathcal H_{a,b}(t-k,\ldots)\right\},
 \label{harm:top}\\
 \mathcal H_{a,b}(\ldots,s,-m)
 &=\frac1q\left\{\sum_{k=0}^q\binom qk B_{q-k}
                  \mathcal H_{a,b}(\ldots,s-k)
                     -B_q(b)\mathcal H_{a,b}(\ldots,s)\right\}.
 \label{harm:bottom}
\end{align}
The singleton rule is
\begin{equation}
             h_{a,b}(-m)=\frac{B_q(a)-B_q(b)}q.
 \label{harm:singleton}
\end{equation}
These are identities of entire functions in every remaining index.
\end{theorem}

\begin{proof}
The classical finite power-sum identity \cite{DLMF} gives
\begin{equation}
 \sum_{z<x<y}x^m=\frac{B_q(y)-B_q(z+1)}q,
 \label{harm:finite-gap}
\end{equation}
when the variables belong to a common unit lattice and the sum has the
indicated strict endpoints. The upper and lower boundary variants
replace $y$ by $a$ and $z+1$ by $b$, respectively. This explains the
formulas for finite sums, but agreement at integer endpoint differences
alone would not justify their analytic continuation. We instead prove
the identities directly in \eqref{harm:normalization}.

For $Z_a$, summing an interior or final nonpositive index by
\eqref{harm:finite-gap} proves \eqref{harm:interior} and
\eqref{harm:bottom}, with $\mathcal H_{a,b}$ replaced by $Z_a$ and
the final lower endpoint $b$ replaced by $a$. One first chooses the
unfixed indices in a domain of absolute convergence and then continues
meromorphically. For a first index use
\[
 \sum_{n_1>n_2}(n_1+a)^{-z}
                   =\zeta(z,n_2+a+1),\qquad
 \zeta(-m,x+1)=-\frac{B_q(x+1)}q.
\]
Taking the remaining indices sufficiently far into their convergence
domain justifies continuation of $z$ to $-m$ inside this expression.
Consequently the top-index identity for $Z_a$ is
\begin{equation}
 Z_a(-m,t,\ldots)
   =-\frac1q\sum_{k=0}^q\binom qk B_{q-k}(1)
                               Z_a(t-k,\ldots).
 \label{harm:Ztop}
\end{equation}
All these statements concern generic values on the indicated spectral
hyperplane, followed by meromorphic continuation; they assign no
directional value to an individual singular $Z_a$ at an intersection.

Now induct on the depth in \eqref{harm:normalization}. For the top
formula substitute \eqref{harm:Ztop} into every prefix of length at
least two. The prefix of length one contributes
$-Z_a(-m)\mathcal H_{a,b}(t,\ldots)
 =B_q(a)\mathcal H_{a,b}(t,\ldots)/q$.
Each remaining group is exactly a shorter convolution identity, giving
\eqref{harm:top}. The bottom formula follows in the same way. Its
two exceptional terms are the last prefix and the singleton suffix;
the factors involving $B_q(a)$ cancel, leaving $B_q(b)$.

For clarity, consider an interior letter and write its nonempty words
on the two sides as $L$ and $R$. In the deconcatenation sum, cuts
strictly before $L$'s last letter use the induction hypothesis on the
suffix; cuts after the first letter of $R$ use the interior $Z_a$
identity on the prefix. The two adjacent cuts are
\[
 Z_a(L)\mathcal H_{a,b}(-m,R),\qquad
 Z_a(L,-m)\mathcal H_{a,b}(R).
\]
After the top and bottom substitutions their exceptional terms are
$B_q(a)Z_a(L)\mathcal H_{a,b}(R)/q$ and its negative.
They cancel. The remaining left-shifted and right-shifted terms are
precisely the recursions for the two shorter words in
\eqref{harm:interior}. This proves that formula on generic remaining
indices. Entireness of $\mathcal H_{a,b}$ then extends all the
identities to every remaining index, including coincident resonances.
\end{proof}

For $m=0$, the interior rule takes the particularly transparent form
\begin{equation}
 \mathcal H_{a,b}(\ldots,s,0,t,\ldots)
 =\mathcal H_{a,b}(\ldots,s-1,t,\ldots)
  -\mathcal H_{a,b}(\ldots,s,t-1,\ldots)
  -\mathcal H_{a,b}(\ldots,s,t,\ldots).
 \label{harm:zero}
\end{equation}
The last term records the strict gap's missing endpoint. Changing
$B_1(1)$ to $B_1$ incorrectly deletes it.

\subsection{A canonical polynomial normal form}

Suppose a word contains $k$ free indices, in their original order
$s_1,\ldots,s_k$, and fixed indices $-m_1,\ldots,-m_e$, where
$m_j\ge0$. Set $M=\sum_{j=1}^e(m_j+1)$.

\Needspace{13\baselineskip}
\begin{theorem}[Finite reduction at arbitrary depth]
\label{harm:normalform}
There are explicitly computable polynomials $C_{\boldsymbol\nu}(a,b)
\in\mathbb Q[a,b]$ such that
\begin{equation}
 \boxed{\displaystyle
 \mathcal H_{a,b}(w)=
  \sum_{\substack{\boldsymbol\nu\in\mathbb N_0^k\\
                    |\boldsymbol\nu|\le M}}
  C_{\boldsymbol\nu}(a,b)
  \mathcal H_{a,b}(s_1-\nu_1,\ldots,s_k-\nu_k),\qquad
 \deg C_{\boldsymbol\nu}\le M-|\boldsymbol\nu|.}
 \label{harm:normalform-eq}
\end{equation}
For $k=0$ the result is a polynomial multiplying the empty word.
Every elimination order in Theorem~\ref{harm:local} gives the same
coefficient polynomials when the free indices are kept formally
distinct. The bound concerns this explicitly specified finite normal
form and does not assert a minimal arithmetic depth.
\end{theorem}

\begin{proof}
Every local rule deletes one fixed index. A deleted index $-m$ has
budget $q=m+1$: it either produces an endpoint polynomial of degree
at most $q$, or shifts an adjacent index by $k\le q$. If the adjacent
index is also fixed, its later budget increases by $k$; if it is free,
its final shift increases by $k$. Induction therefore proves termination
and the degree bound in \eqref{harm:normalform-eq}.

Here is a canonical description of the coefficients. For a fixed
nonpositive block $v$, let
\[
 \Pi_v(A,B)=\mathcal H_{A,B}(v),\qquad \Pi_{\varnothing}(A,B)=1.
\]
The first part proves that these are polynomials. Write
$w=(v_0,s_1,v_1,\ldots,s_k,v_k)$ and expand the polynomial
\begin{equation}
 G_w(\boldsymbol y;a,b)=
   \Pi_{v_0}(a,y_1+1)
   \prod_{j=1}^{k-1}\Pi_{v_j}(y_j,y_{j+1}+1)
   \Pi_{v_k}(y_k,b)
   =\sum_{\boldsymbol\nu}C_{\boldsymbol\nu}(a,b)
                                 \boldsymbol y^{\boldsymbol\nu}.
 \label{harm:gaps}
\end{equation}
For integer endpoint differences, fix the free summation variables
$y_1>\cdots>y_k$. The sums in the intervening gaps are independent,
and their product is exactly \eqref{harm:gaps}.

To identify these coefficients with those of any already proved local
reduction, compare the two finite sums for all independent $s_j$.
Distinct ordered tuples of positive lattice points give linearly
independent functions $\prod_j y_j^{-s_j}$: this is elementary finite
linear independence of exponential functions, applied successively
in the $s_j$. Thus the difference of the two coefficient polynomials
vanishes at every lattice configuration
$a=b+N$, $y_j=b+n_j$, $N>n_1>\cdots>n_k\ge0$, with $b>0$.
These configurations are Zariski dense. Indeed a polynomial vanishing
on all integer tuples satisfying strict inequalities vanishes by
successive univariate polynomial uniqueness, and $b$ ranges over an
interval. Hence the coefficient polynomials agree identically.
This is polynomial uniqueness after the analytic identities have
already been proved; it is not uniqueness of a holomorphic interpolation
from integer endpoint values.
\end{proof}

In the one-free-index case, every depth reduces to the finite sum
\begin{equation}
           \mathcal H_{a,b}(v_0,s,v_1)
              =\sum_{\nu=0}^M C_\nu(a,b)h_{a,b}(s-\nu).
 \label{harm:onefree}
\end{equation}
All derivatives in the free index commute with this identity. At
$s-\nu=1$ use the finite value
\begin{equation}
 h_{a,b}^{(r)}(1)=(-1)^r\{\gamma_r(b)-\gamma_r(a)\},
 \qquad
 h_{a,b}'(0)=\log\Gamma(b)-\log\Gamma(a).
 \label{harm:jets}
\end{equation}
At positive integer arguments $>1$, the undifferentiated differences
are polygamma differences. These standard Hurwitz specializations are
recorded in \cite{DLMF}. Differentiation of a \emph{fixed}
index $-m$ is a different problem and is not licensed by a formula
proved only on that spectral hyperplane.

Two concise examples are
\begin{align}
 \mathcal H_{a,b}(0,s,0)
   &=-h_{a,b}(s-2)+(a+b-1)h_{a,b}(s-1)
                                  -b(a-1)h_{a,b}(s),
 \label{harm:examplezero}\\
 \mathcal H_{a,b}(-1,s,-1)
   &=\frac14\bigl\{-h_{a,b}(s-4)
       +(A+B+1)h_{a,b}(s-2)\notag\\
   &\hspace{32mm}+(B-A)h_{a,b}(s-1)-ABh_{a,b}(s)\bigr\}.
 \label{harm:exampleone}
\end{align}
Here $A=a(a-1)$ and $B=b(b-1)$ in the second identity.
The absent shift $s-3$ in the second identity is an exact cancellation.
It follows by expanding
$(A-y^2-y)(y^2-y-B)/4$. For example, its derivative at $s=2$ is
\begin{align}
 \left.\partial_s\mathcal H_{a,b}(-1,s,-1)\right|_{s=2}
  =\frac14\bigl\{&-\zeta'(-2,b)+\zeta'(-2,a)\notag\\
       &+(A+B+1)\log\tfrac{\Gamma(b)}{\Gamma(a)}
        +(B-A)\{\gamma_1(a)-\gamma_1(b)\}\notag\\
       &-AB\{\zeta'(2,b)-\zeta'(2,a)\}\bigr\}.
 \label{harm:examplejet}
\end{align}

\subsection{Summing arbitrarily many zero indices}

For a block of $p$ zero indices, polynomial continuation of the finite
count gives
\begin{equation}
                    \Pi_{\{0\}^p}(A,B)=\binom{A-B}{p}.
 \label{harm:zeroblock}
\end{equation}
Consequently \eqref{harm:gaps} for a single free index uses the polynomial
$\binom{a-y-1}{p}\binom{y-b}{q}$. All such polynomials can be summed
in one analytic formula.

\begin{theorem}[Lerch generator for zero decorations]
\label{harm:Lerch}
Let $\Re a,\Re b>0$. There is a function
$\mathcal L_{a,b}(s;t)$ entire in $s$ and holomorphic for $|t|<2\pi$
such that, when $\Re t<0$,
\begin{equation}
 \mathcal L_{a,b}(s;t)=
    e^{bt}\Phi(e^t,s,b)-e^{at}\Phi(e^t,s,a).
 \label{harm:Lerchdifference}
\end{equation}
Here $\Phi(z,s,c)=\sum_{n\ge0}z^n(n+c)^{-s}$ for $|z|<1$, with
principal powers of $n+c$. The combination in
\eqref{harm:Lerchdifference}, not its separate summands, is continued
through $t=0$. It satisfies
\begin{equation}
 \mathcal L_{a,b}(s;0)=h_{a,b}(s),\qquad
 \partial_t\mathcal L_{a,b}(s;t)=\mathcal L_{a,b}(s-1;t),\qquad
 \mathcal L_{a,b}(s;t)=\sum_{n\ge0}h_{a,b}(s-n)\frac{t^n}{n!}.
 \label{harm:LerchTaylor}
\end{equation}
The last series converges locally normally for $|t|<2\pi$, including
all fixed spectral derivatives. For $|u|,|v|<1/2$ one has the convergent
identity
\begin{equation}
 \boxed{\displaystyle
 \sum_{p,q\ge0}\mathcal H_{a,b}(\{0\}^p,s,\{0\}^q)u^pv^q
  =(1+u)^{a-1}(1+v)^{-b}
    \mathcal L_{a,b}\bigl(s;\log(1+v)-\log(1+u)\bigr).}
 \label{harm:zerogenerator}
\end{equation}
Powers and logarithms on the right are their germs at $u=v=0$.
\end{theorem}

\begin{proof}
Define
\begin{equation}
 K_{a,b}(y)=\frac{e^{-by}-e^{-ay}}{1-e^{-y}}.
 \label{harm:K}
\end{equation}
The apparent pole at $y=0$ is removable, with value $a-b$.
Its only other possible finite poles are $2\pi i\mathbb Z\setminus\{0\}$.
Thus for $|t|<2\pi$ the function $K_{a,b}(x-t)$ is analytic near
every $x\ge0$ and decreases exponentially as $x\to+\infty$,
locally uniformly in $t,a,b$. Initially for $\Re s>0$ put
\begin{equation}
 \mathcal L_{a,b}(s;t)=\frac1{\Gamma(s)}
                \int_0^\infty x^{s-1}K_{a,b}(x-t)\,dx.
 \label{harm:LerchMellin}
\end{equation}
For any $N\ge1$, subtract the Taylor polynomial of $K_{a,b}(x-t)$
at $x=0$ in the integral from $0$ to $1$, and add
\[
 \frac1{\Gamma(s)}\sum_{j=0}^{N-1}
             \frac{K_{a,b}^{(j)}(-t)}{j!(s+j)}.
\]
The subtracted integral is holomorphic for $\Re s>-N$. The displayed
possible poles are cancelled by the simple zeros of $1/\Gamma(s)$.
Since $N$ is arbitrary, this proves joint entireness in $s$ and
holomorphy for $|t|<2\pi$, with locally normal spectral derivatives.

For $\Re t<0$, geometric expansion in \eqref{harm:LerchMellin}
and the Gamma integral give \eqref{harm:Lerchdifference}, first for
$\Re s>0$ and then for all $s$. At $t=0$ the ordinary Hurwitz
integral gives $h_{a,b}(s)$, initially for $\Re s>1$, and then by
continuation. Differentiating the absolutely convergent Lerch series
for $\Re t<0$ proves the order-lowering identity in
\eqref{harm:LerchTaylor}; continuation and Taylor's theorem give the
remaining assertions.

The right side of \eqref{harm:zerogenerator} is holomorphic for
$|u|,|v|<1/2$, because the absolute value of the difference of its
logarithms is less than $2\log2<2\pi$. To identify a coefficient,
formally replace $h_{a,b}(s-n)$ in \eqref{harm:LerchTaylor} by
$y^n$. Its right side becomes $e^{yt}$. Including the prefactors in
\eqref{harm:zerogenerator} gives
\[
                 (1+u)^{a-y-1}(1+v)^{y-b}.
\]
Its coefficient of $u^pv^q$ is precisely
$\binom{a-y-1}{p}\binom{y-b}{q}$. At each fixed bidegree this is a
finite polynomial calculation. The canonical gap formula then identifies
it with $\mathcal H_{a,b}(\{0\}^p,s,\{0\}^q)$, proving the
identity and the convergence of its left side.
\end{proof}

The branch cancellation can also be read as a Gamma separation:
$\Gamma(s)\mathcal L_{a,b}(s;t)$ is the ordinary Mellin transform
of the explicit kernel \eqref{harm:K}; its apparent spectral poles at
nonpositive integers are entirely accounted for by $\Gamma(s)$.
This proves a generating formula for the specified one-free-order
family. It does not settle a minimality question for several arbitrary
independent nonzero indices.

Setting $u=v$ gives the exact insertion identity
\begin{equation}
 \sum_{p+q=r}\mathcal H_{a,b}(\{0\}^p,s,\{0\}^q)
                  =\binom{a-b-1}{r}h_{a,b}(s).
 \label{harm:diagonalzero}
\end{equation}
The generating function also has normalized antiderivatives:
\begin{equation}
 \int_0^t\mathcal L_{a,b}(s;\tau)\,d\tau
            =\mathcal L_{a,b}(s+1;t)-h_{a,b}(s+1).
 \label{harm:Lerchprimitive}
\end{equation}
At a nonpositive order,
\begin{equation}
 \mathcal L_{a,b}(-m;t)
      =\partial_t^m\frac{e^{bt}-e^{at}}{1-e^t},
 \qquad m\ge0,
 \label{harm:Lerchnegative}
\end{equation}
where $t=0$ is removable. At $a=b+N$, $N\ge0$ integer,
$\mathcal L_{a,b}(s;t)=\sum_{n=0}^{N-1}(n+b)^{-s}e^{(n+b)t}$;
the double generating function becomes a polynomial of total degree at
most $N-1$ when $N\ge1$ and is zero when $N=0$.

\subsection{Polynomial endpoint primitives and Stieltjes moments}

The polynomial Hurwitz integration mechanism used here is classical:
Espinosa and Moll give a general integration-by-parts formula and its
polynomial specialization \cite[Theorem~2.1 and Example~2.2]{harm:EM}.
The canonical manuscript \cite{source:manuscript}, in the Moments section
of \src{chapters/07-integration.tex}, already states the unit-interval
Hurwitz transform, gives the first Stieltjes moment, and identifies the complete
homogeneous coefficients for every Stieltjes order. We state the fully
indexed arbitrary-endpoint version, with the normalization required by
\eqref{harm:onefree}. This is an explicit extension and restatement of
the existing moment mechanism, and is not counted as a new general
Stieltjes-moment principle.

For $d\ge0$, write $\Delta=a-b$ and define
\begin{equation}
 P_{d;a,b}(s)=\int_b^a(x-b)^d h_{x,b}(s)\,dx.
 \label{harm:Pdef}
\end{equation}
Paths lie in the right half-plane. The integral is path independent and
entire in $s$. Initially away from singularities of the individual
displayed terms,
\begin{align}
 P_{d;a,b}(s)
  ={}&\frac{\Delta^{d+1}}{d+1}\zeta(s,b)\notag\\
   &+\sum_{j=0}^d\frac{d!}{(d-j)!}
     \frac{\Delta^{d-j}\zeta(s-j-1,a)
              -\mathbf1_{j=d}\zeta(s-j-1,b)}
          {(s-1)(s-2)\cdots(s-j-1)}.
 \label{harm:Pfinite}
\end{align}
This follows by repeated integration by parts using
$\partial_x\zeta(s-1,x)=-(s-1)\zeta(s,x)$.
The definite integral proves that every apparent pole in the combined
right side is removable. Expanding a coefficient polynomial
$C_\nu(x,b)$ in powers of $x-b$ and using \eqref{harm:Pfinite}
therefore integrates every one-free-order word and every one of its
free-order derivatives in a finite Hurwitz basis, with the constant
fixed at $a=b$.

To state all resonance coefficients, put
\begin{equation}
 c_{j,n}=[z^n]\prod_{\ell=1}^j(1-z/\ell)^{-1},
 \qquad c_{0,n}=\mathbf1_{n=0}.
 \label{harm:cn}
\end{equation}
These are finite generalized harmonic expressions:
\[
 \sum_{n\ge0}c_{j,n}z^n
     =\exp\left(\sum_{r\ge1}\frac{H_j^{(r)}}r z^r\right),
 \qquad H_j^{(r)}=\sum_{\ell=1}^j\ell^{-r}.
\]
For example $c_{j,0}=1$, $c_{j,1}=H_j$ and
$c_{j,2}=(H_j^2+H_j^{(2)})/2$. Equivalently, $c_{j,n}$ is the
complete exponential Bell polynomial in
$H_j,1!H_j^{(2)},\ldots,(n-1)!H_j^{(n)}$, divided by $n!$.

\begin{corollary}[Explicit arbitrary-endpoint form of the moment calculus]
\label{harm:Stieltjesmoments}
For $d,m\ge0$ and $\Re a,\Re b>0$,
\begin{align}
 \int_b^a(x-b)^d\gamma_m(x)\,dx
   ={}&(-1)^{m+1}m!
        \sum_{j=0}^d(-1)^j\binom dj
        \sum_{r=0}^{m+1}\frac{c_{j,m+1-r}}{r!}\notag\\
      &\hspace{7mm}\cdot
       \left\{\Delta^{d-j}\zeta^{(r)}(-j,a)
                      -\mathbf1_{j=d}\zeta^{(r)}(-j,b)\right\}.
 \label{harm:Stieltjesfinite}
\end{align}
Here derivatives of $\zeta$ are with respect to its first argument.
\end{corollary}

\begin{proof}
Set $s=1+z$ in \eqref{harm:Pfinite}. The denominator in its $j$th
term has expansion
\[
 \frac1{z(z-1)\cdots(z-j)}
       =\frac{(-1)^j}{j!z}\sum_{n\ge0}c_{j,n}z^n.
\]
Use
$\zeta(1+z,b)=z^{-1}+\sum_{m\ge0}(-1)^m\gamma_m(b)z^m/m!$
and the ordinary Taylor expansion of each $\zeta(z-j,a)$.
The coefficient of $z^{-1}$ cancels, as it must by entireness of
$P_{d;a,b}$. The coefficient of $z^m$, multiplied by $m!$, equals
\[
 (-1)^m\frac{\Delta^{d+1}}{d+1}\gamma_m(b)
 +m!\sum_{j=0}^d(-1)^j\binom dj
    \sum_{r=0}^{m+1}\frac{c_{j,m+1-r}}{r!}
       \left\{\Delta^{d-j}\zeta^{(r)}(-j,a)
                         -\mathbf1_{j=d}\zeta^{(r)}(-j,b)\right\}.
\]
On the integral side, \eqref{harm:jets} gives
$(-1)^m\{\Delta^{d+1}\gamma_m(b)/(d+1)
 -\int_b^a(x-b)^d\gamma_m(x)\,dx\}$.
Cancel the common endpoint term and solve for the integral.
All coefficient extractions are locally normal on compact endpoint
sets in the right half-plane.
\end{proof}

For $d=0$ this recovers the familiar normalized primitive
\begin{equation}
 \int_b^a\gamma_m(x)\,dx
       =\frac{(-1)^{m+1}}{m+1}
          \{\zeta^{(m+1)}(0,a)-\zeta^{(m+1)}(0,b)\}.
 \label{harm:d0}
\end{equation}
For $d=1$, since $c_{1,n}=1$ for every $n$,
\begin{align}
 \int_b^a(x-b)\gamma_m(x)\,dx
  =(-1)^{m+1}m!\left\{
       \frac{\Delta\zeta^{(m+1)}(0,a)}{(m+1)!}
       -\sum_{r=0}^{m+1}
          \frac{\zeta^{(r)}(-1,a)-\zeta^{(r)}(-1,b)}{r!}
                         \right\}.
 \label{harm:d1}
\end{align}
In particular $m=0$ reduces to an integral of the digamma function,
and Lerch's identity converts the $\zeta'(0,a)$ term to
$\log\Gamma(a)-\tfrac12\log(2\pi)$.

\subsection{Exact scope and further questions}

The identities above add arbitrary-depth finite reductions and a
convergent all-zero-decoration generator to the previously established
depth-two calculus. They do not evaluate nonsymmetric positive-index
coordinates such as the incoming harmonic Stieltjes parameter $\eta$.
No statement about the unresolved $S_6$ or $S_8$ reductions is made.

Natural next questions are the following.
\begin{enumerate}
\item Construct and prove a jointly convergent analogue of
\eqref{harm:zerogenerator} with two or more independently moving free
indices and arbitrary zero blocks between them. The finite coefficient
polynomials \eqref{harm:gaps} are already explicit; the remaining task
is a useful analytic kernel and a precise simultaneous branch domain.
\item Differentiate an eliminated index before specializing it to
$-m$. Determine a finite set of additional ordered coordinates needed
at the first such derivative. The present identities only differentiate
the free indices, so this question asks for a real extension.
\item For rational endpoints, combine \eqref{harm:Stieltjesfinite}
with finite character projections to give algorithms for polynomial
moments in Dirichlet $L$-derivatives, retaining every conductor logarithm.
\item Formalize the Bernoulli deletion rules, the polynomial gap
normal form, and the finite coefficient extraction. Their analytic
interface is the already established entire interpolant and the
explicit normalized Mellin continuation in \eqref{harm:LerchMellin}.
\end{enumerate}

% END sections/04_harmonic_reduction.tex


% BEGIN sections/05_tornheim.tex
% Independent research contribution for the parent article.
% Required packages: amsmath,amssymb,amsthm; ordinary theorem environments.
% All labels and bibliography keys have qtr: prefix.
\section{The complete quartic Tornheim layer and the formal jets at every order}
\label{qtr:section}

The incoming \emph{Independent Orders and Exact Reductions} report
reduces every cubic directional derivative of the Tornheim function to
two scalar coordinates and asks for its quartic layer
\cite{qtr:Independent}. We determine that layer explicitly. Three
coordinates suffice: the diagonal fourth derivative, one convergent
logarithmic Gamma moment, and one absolutely convergent integral of
ordinary polylogarithmic order derivatives. We also determine the exact
dimension of the homogeneous jet space allowed by the source's symmetry,
axis, and Euler-plane identities. These are statements about a
specified formal relation space, not about arithmetic independence of
the constants.

\subsection{The analytic normalization and its exact restrictions}

Write
\[
 T(A,B,C)=\sum_{m,n\ge1}m^{-A}n^{-B}(m+n)^{-C},\qquad
 L=\log(2\pi),\quad z_j=\zeta^{(j)}(0),\quad
 y_j=\zeta^{(j)}(-1).
\]
The series is initially used only in its domain of absolute
convergence; elsewhere \(T\) denotes its meromorphic continuation.
The origin is not a joint regular point. For a direction
\((a,b,c)\) satisfying \((a+c)(b+c)\ne0\), put
\[
 W_{a,b,c}(t)=T(at,bt,ct),
\]
where the restriction is formed before continuation to \(t=0\).

The source constructs the unique symmetric holomorphic germ
\(R(A,B,C)=R(B,A,C)\) such that
\begin{equation}\label{qtr:polar}
 T(A,B,C)=\zeta(A)\zeta(B)
 -\frac{C\zeta(B-1)}{A+C}
 -\frac{C\zeta(A-1)}{B+C}+C R(A,B,C).
\end{equation}
For completeness, Jonqui\`ere's expansion \cite[25.12.12]{DLMF}
gives this by subtracting the following six local terms from the
Mellin kernel:
\begin{align}
 \mathcal S(A,B;x)={}&
 \Gamma(1-A)\Gamma(1-B)x^{A+B-2}
 +\Gamma(1-A)\zeta(B)x^{A-1}
 +\Gamma(1-B)\zeta(A)x^{B-1}\notag\\
 &-\Gamma(1-A)\zeta(B-1)x^A
 -\Gamma(1-B)\zeta(A-1)x^B+\zeta(A)\zeta(B).
 \label{qtr:S}
\end{align}
The subtracted remainder and every fixed parameter derivative are
normally integrable after multiplication by \(x^{C-1}\) in a
sufficiently small common polydisc. Before removing the Gamma
factors the polar terms contain
\(\Gamma(1-A)/\Gamma(1+C)\). Their difference from \(1\) is
divisible by \(A+C\), since it vanishes when \(C=-A\).
This proves that the difference between the resulting polar model
and \eqref{qtr:polar}, divided by \(C\), is holomorphic.

We will use precisely two functional restrictions of this germ:
\begin{align}
 C R(0,0,C)&=\zeta(C-1)-\zeta(C)-\frac5{12},
 \label{qtr:axis}\\
 C R(A,0,C)+A R(C,0,A)
 &=\zeta(A)\zeta(C)-\zeta(A+C)
 +\frac{\zeta(A)+\zeta(C)}2\notag\\
 &\quad+\zeta(-1)+\zeta(A-1)+\zeta(C-1).
 \label{qtr:Euler}
\end{align}
The first follows by counting pairs with prescribed sum.
The second follows from the Euler stuffle formula for
\(T(A,0,C)+T(C,0,A)\). Both are holomorphic identities near the
origin. They were proved in the incoming report; their use here
does not constitute a new proof of the cubic theorem.

\subsection{All quartic directions in three explicit coordinates}

All partial derivatives of \(R\) below are taken at the origin.
Set
\[
 \Lambda=W_{1,1,1}^{(4)}(0),\qquad
 U=R_{AAA},\qquad \Xi=R_{AAB}.
\]
The notation \(\Xi\) distinguishes this new coordinate from the
several unrelated uses of \(\kappa\) in the source articles.

\begin{lemma}[The degree-three remainder constraints]
\label{qtr:constraints}
Write
\begin{align*}
 u&=R_{AAA}=R_{BBB},& v&=R_{AAB}=R_{ABB},&
 w&=R_{AAC}=R_{BBC},\\
 x&=R_{ABC},& y&=R_{ACC}=R_{BCC},& z&=R_{CCC}.
\end{align*}
Then
\begin{align}
 z&=\frac{y_4-z_4}{4},&
 w&=\frac{z_2^2-z_4}{4},&
 u+3y&=-\frac L2z_3-z_4,\label{qtr:jetconstraints}\\
 v+x&=\frac{\Lambda+8Lz_3-12z_2^2+16z_4}{24}.
 \label{qtr:diagonalconstraint}
\end{align}
\end{lemma}

\begin{proof}
The coefficient of \(C^4\) in \eqref{qtr:axis} is \(z/6\),
and its right-hand side is \((y_4-z_4)/24\). In
\eqref{qtr:Euler}, the coefficients of \(A^2C^2\) and
\(A^3C\) are respectively \(w\) and \(u/6+y/2\).
The corresponding right-hand coefficients are
\((z_2^2-z_4)/4\) and
\((z_1z_3-z_4)/6\), where \(z_1=-L/2\).
This proves \eqref{qtr:jetconstraints}.

On the diagonal, \eqref{qtr:polar} becomes
\[
 T(t,t,t)=\zeta(t)^2-\zeta(t-1)+tR(t,t,t).
\]
Its fourth derivative is
\[
 \Lambda=-z_4-4Lz_3+6z_2^2-y_4
             +8u+24v+24w+24x+24y+4z.
\]
Substituting \eqref{qtr:jetconstraints} leaves
\(\Lambda=-8Lz_3+12z_2^2-16z_4+24(v+x)\),
as required.
\end{proof}

\begin{theorem}[Every quartic Tornheim ray]\label{qtr:quartic}
For \((a+c)(b+c)\ne0\), the restricted function \(W_{a,b,c}\)
is holomorphic at zero and
\begin{align}
 W_{a,b,c}^{(4)}(0)={}&
 abc^2\Lambda
 +4c\bigl(a^3+b^3-c^2(a+b)\bigr)U
 +12abc(a+b-2c)\Xi\notag\\
 &-2\bigl(ab(a^2+b^2-4c^2)+c^3(a+b)\bigr)Lz_3\notag\\
 &+3\bigl(2a^2b^2+c^2(a^2+b^2-4ab)\bigr)z_2^2\notag\\
 &+\left(-\frac{a^4+b^4}{2}+16abc^2
        -3c^2(a^2+b^2)-4c^3(a+b)-c^4\right)z_4\notag\\
 &+\left(c^4-\frac{cb^4}{a+c}-\frac{ca^4}{b+c}\right)y_4.
 \label{qtr:quarticformula}
\end{align}
The scalar \(U\) is given by the convergent Gamma series
\eqref{qtr:U}; \(\Xi\) is given by the convergent integral
\eqref{qtr:Xi}. No unevaluated partial derivative of \(T\)
appears in these two coordinates.
\end{theorem}

\begin{proof}
On an admissible ray the two rational polar terms in
\eqref{qtr:polar} become constant multiples of holomorphic
one-variable zeta functions, proving the asserted holomorphy.
The degree-three homogeneous part of \(R\) is
\[
 \frac{u}{6}(A^3+B^3)+\frac{v}{2}AB(A+B)
 +\frac{w}{2}C(A^2+B^2)+xABC
 +\frac{y}{2}C^2(A+B)+\frac{z}{6}C^3.
\]
Therefore the fourth derivative of \eqref{qtr:polar} is
\begin{align*}
 &-\frac{a^4+b^4}{2}z_4
 -2ab(a^2+b^2)Lz_3+6a^2b^2z_2^2
 -c\left(\frac{b^4}{a+c}+\frac{a^4}{b+c}\right)y_4\\
 &\qquad+4c\left(u(a^3+b^3)+3vab(a+b)
 +3wc(a^2+b^2)+6xabc+3yc^2(a+b)+zc^3\right).
\end{align*}
Put \(u=U,v=\Xi\) and substitute
\eqref{qtr:jetconstraints}--\eqref{qtr:diagonalconstraint}.
Collecting the five ordinary zeta combinations and the three
retained coordinates gives \eqref{qtr:quarticformula}.
\end{proof}

\paragraph{A quartic cyclic identity with no additional coordinates.}
Let \(a+b,a+c,b+c\) all be nonzero and set
\(e_1=a+b+c\), \(e_2=ab+bc+ca\), \(e_3=abc\).
Then
\begin{align}
 \sum_{\mathrm{cyc}}W_{a,b,c}^{(4)}(0)
 ={}&e_1e_3\Lambda
 +(-4e_1^2e_2+8e_2^2+12e_1e_3)Lz_3\notag\\
 &+(12e_2^2-36e_1e_3)z_2^2\notag\\
 &+(-2e_1^4+4e_1^2e_2-2e_2^2+24e_1e_3)z_4.
 \label{qtr:cyclicquartic}
\end{align}
Indeed, the cyclic coefficients of \(U\) and \(\Xi\) in
\eqref{qtr:quarticformula} vanish identically. Pairing the two
terms with each denominator cancels every \(y_4\) term.
The remaining polynomial coefficients give the displayed formula.
For example,
\begin{equation}\label{qtr:example}
 W_{1,2,3}^{(4)}(0)+W_{2,3,1}^{(4)}(0)
 +W_{3,1,2}^{(4)}(0)-36\Lambda
 =-184Lz_3+156z_2^2-386z_4.
\end{equation}
The Euler diagonal supplies an independent elementary check:
\[
 W_{1,0,1}^{(4)}(0)
 =-2Lz_3+3z_2^2-\frac{17}{2}z_4,
\]
which follows directly from
\(T(t,0,t)=(\zeta(t)^2-\zeta(2t))/2\).

More generally, the exact coordinate map
\[
 (a,b,c)\longmapsto
 \bigl(abc^2,\;
 4c(a^3+b^3-c^2(a+b)),\;
 12abc(a+b-2c)\bigr)
\]
gives a sufficient algebraic criterion for finite linear
combinations of quartic rays to reduce entirely to the ordinary
zeta derivatives displayed in \eqref{qtr:quarticformula}.
This criterion is exact in the stated three-coordinate normal
form. Additional arithmetic relations among the coordinates
could yield further reductions.

\subsection{Convergent representatives for the two additional coordinates}

Define the Stirling remainder
\[
 \rho(n)=\log\Gamma(n+1)-\left(n+\frac12\right)\log n+n
                 -\frac L2-\frac1{12n}.
\]
The source's boundary-germ theorem states, for every \(j\ge0\),
\[
 \partial_A^jR(0,0,0)
 =(-1)^j\sum_{n\ge1}(\log n)^j\rho(n)
 -y_{j+1}-\frac12z_{j+1}-y_j
 +\frac{(-1)^j\gamma_j}{12}.
\]
Consequently the coordinate needed here is
\begin{equation}\label{qtr:U}
 U=-\sum_{n\ge1}(\log n)^3\rho(n)
                  -y_4-\frac12z_4-y_3-\frac{\gamma_3}{12}.
\end{equation}
This is a specialization of an existing result, not a separate
novel Gamma-moment identity. Absolute convergence follows from
\(\rho(n)=O(n^{-3})\).

For the genuinely mixed coordinate let
\[
 f_j(x)=\left.\partial_s^j\operatorname{Li}_s(e^{-x})\right|_{s=0},
 \qquad h(x)=\log x+\gamma,\qquad j=1,2.
\]
For fixed \(x>0\), the defining series are absolutely convergent:
\[
 f_1(x)=-\sum_{n\ge1}(\log n)e^{-nx},\qquad
 f_2(x)=\sum_{n\ge1}(\log n)^2e^{-nx}.
\]
Put
\begin{align}
 P(x)={}&\frac{h(x)^3+\zeta(2)h(x)}{x^2}
 +\frac{z_1(h(x)^2+\zeta(2))+z_2h(x)}{x}\notag\\
 &-y_1(h(x)^2+\zeta(2))-y_2h(x)+z_1z_2,
 \label{qtr:P}
\end{align}
and define the ordinary absolutely convergent integral
\begin{equation}\label{qtr:I}
 I=\int_0^1\frac{f_2(x)f_1(x)-P(x)}{x}\,dx
                  +\int_1^\infty\frac{f_2(x)f_1(x)}{x}\,dx.
\end{equation}

\begin{proposition}[Explicit mixed polylogarithmic moment]
\label{qtr:Xiprop}
Let \(g=\gamma\), \(h_2=\zeta(2)\), \(h_3=\zeta(3)\), and set
\begin{align}
 E={}&-\frac{g^3+gh_2}{2}
       -\frac{3g^2+h_2}{4}-\frac{3g}{4}-\frac38\notag\\
 &-(g^2+h_2+2g+2)z_1-(g+1)z_2+gz_1z_2\notag\\
 &-\frac{g^3+3gh_2+2h_3}{3}y_1
       -\frac{g^2+h_2}{2}y_2.
 \label{qtr:E}
\end{align}
Then
\begin{equation}\label{qtr:Xi}
 \Xi=I+E.
\end{equation}
\end{proposition}

\begin{proof}
Define the normally convergent subtracted Mellin integral
\[
 I(A,B)=\int_0^1
 \frac{\operatorname{Li}_A(e^{-x})\operatorname{Li}_B(e^{-x})
                    -\mathcal S(A,B;x)}{x}\,dx
 +\int_1^\infty
 \frac{\operatorname{Li}_A(e^{-x})\operatorname{Li}_B(e^{-x})}{x}\,dx.
\]
The \(C=0\) specialization of the proof of \eqref{qtr:polar} gives
the holomorphic identity
\begin{align}
 R(A,B,0)={}&I(A,B)
 +\frac{\Gamma(1-A)\Gamma(1-B)}{A+B-2}
 +\frac{\Gamma(1-A)\zeta(B)}{A-1}
 +\frac{\Gamma(1-B)\zeta(A)}{B-1}\notag\\
 &+\gamma\zeta(A)\zeta(B)
 -\frac{\Gamma(1-A)-1}{A}\zeta(B-1)
 -\frac{\Gamma(1-B)-1}{B}\zeta(A-1).
 \label{qtr:boundarymixed}
\end{align}
Both apparent quotients at \(A=0\) or \(B=0\) are removable.
Differentiating \(\mathcal S\) twice in \(A\) and once in \(B\)
at zero produces precisely \(P(x)\). Thus
\(I_{AAB}(0,0)=I\).

To see convergence directly, write the order-differentiated
Jonqui\`ere expansion as
\[
 f_j(x)=\frac{g_j(x)}{x}
       +\sum_{k\ge0}\frac{(-1)^k\zeta^{(j)}(-k)}{k!}x^k,
 \qquad
 g_1=h,\quad g_2=h^2+\zeta(2).
\]
The subtraction \(P\) removes the singular product, both
cross terms with \(k=0,1\), and the constant analytic product.
The remaining numerator is
\(O(x(1+|\log x|^2))\) at zero. At infinity
\(f_1f_2=O(e^{-4x})\). This proves ordinary absolute convergence;
normal convergence in a parameter polydisc justifies the
differentiation used above.

Finally,
\[
 \Gamma(1-t)=1+gt+\frac{g^2+h_2}{2}t^2
               +\frac{g^3+3gh_2+2h_3}{6}t^3+O(t^4).
\]
Applying \(\partial_A^2\partial_B\) to the six elementary
terms after \(I(A,B)\) in \eqref{qtr:boundarymixed}, using this
Taylor expansion, gives exactly \(E\). This proves
\eqref{qtr:Xi}.
\end{proof}

The representation \eqref{qtr:Xi} fixes the coordinate using only
ordinary polylogarithms, their order derivatives, Gamma constants,
and ordinary zeta derivatives. It is not an asserted reduction
of \(\Xi\) to a finite algebra of standard constants.
The splitting point \(1\) in \eqref{qtr:I} is part of the formula;
changing it requires changing the explicitly integrated
subtractions.

\subsection{The formal relation space at arbitrary degree}

Here a \emph{formal degree-\(d\) freedom} means the difference of
two homogeneous degree-\(d\) jets satisfying symmetry in \(A,B\),
the same degree of \eqref{qtr:axis}, and the same degree of
\eqref{qtr:Euler}. No other functional relation is imposed.

\begin{theorem}[Exact homogeneous kernel]\label{qtr:kernel}
For \(d\ge2\), every formal degree-\(d\) freedom has a unique
representation
\begin{equation}\label{qtr:kernelform}
 \Delta R_d(A,B,C)=
 A(A-C)S(A,C)+B(B-C)S(B,C)+ABQ(A,B,C),
\end{equation}
where \(S\) is a homogeneous symmetric polynomial of degree
\(d-2\) in two variables and \(Q\) is a homogeneous polynomial of
degree \(d-2\), symmetric in \(A,B\). Therefore the dimension of
this formal freedom is
\begin{equation}\label{qtr:dimension}
 \dim\mathcal K_d
 =\left\lfloor\frac d2\right\rfloor
     +\left\lfloor\frac{d^2}{4}\right\rfloor.
\end{equation}
At degrees \(0\) and \(1\) the freedom is zero.
\end{theorem}

\begin{proof}
Put \(D(A,C)=\Delta R_d(A,0,C)\). The axis condition gives
\(D(0,C)=0\), so \(D=A V(A,C)\). The homogeneous Euler
condition becomes
\[
 AC\bigl(V(A,C)+V(C,A)\bigr)=0.
\]
Hence \(V\) is antisymmetric in \(A,C\), and
\(V(A,C)=(A-C)S(A,C)\), where \(S\) is symmetric.
This proves the asserted boundary value.

Subtract the first two terms of \eqref{qtr:kernelform} from
\(\Delta R_d\). The difference is still symmetric in \(A,B\)
and vanishes when \(A=0\) or \(B=0\). It is therefore divisible
by \(AB\); its quotient is the required symmetric polynomial
\(Q\). This construction determines \(S\) from the boundary
and then determines \(Q\), proving uniqueness. Conversely, a
direct substitution shows that every expression
\eqref{qtr:kernelform} satisfies all three homogeneous
conditions.

The space of symmetric polynomials of degree \(d-2\) in two
variables has dimension \(\lfloor d/2\rfloor\). The space
of degree-\(d-2\) polynomials in three variables symmetric
in the first two has dimension
\[
 \sum_{k=0}^{d-2}
 \left(\left\lfloor\frac{d-2-k}{2}\right\rfloor+1\right)
 =\left\lfloor\frac{d^2}{4}\right\rfloor.
\]
For \(d=0,1\), the divisibility argument forces \(S\) and \(Q\)
to vanish.
\end{proof}

The Tornheim derivative of order \(n\) uses degree \(d=n-1\)
of \(R\). Thus the first formal dimensions, for derivative
orders \(n=1,\ldots,8\), are
\[
 0,\quad0,\quad2,\quad3,\quad6,\quad8,\quad12,\quad15.
\]
In particular the quartic three-coordinate count is exact under
these stated relations. This count does not say that three
arithmetically independent periods are necessary.

\begin{corollary}[All-order cyclic cancellation and its exact image]
\label{qtr:cyclicimage}
For \(d\ge2\), the image of the homogeneous freedom under
\[
 \Delta R_d\longmapsto
 c\Delta R_d(a,b,c)+a\Delta R_d(b,c,a)
                         +b\Delta R_d(c,a,b)
\]
is exactly \(abc\) times the space of symmetric homogeneous
polynomials of degree \(d-2\) in \(a,b,c\). Consequently its
dimension is
\[
 p_{\le3}(d-2)=
 \left\lfloor\frac{(d+1)^2+3}{12}\right\rfloor.
\]
Equivalently, at derivative order \(n\ge3\), cyclic
symmetrization leaves exactly
\(\lfloor(n^2+3)/12\rfloor\) formal coordinates.
\end{corollary}

\begin{proof}
The terms involving \(S\) in \eqref{qtr:kernelform} cancel
pairwise because \(S(a,c)=S(c,a)\). The remaining expression is
\[
 abc\bigl(Q(a,b,c)+Q(b,c,a)+Q(c,a,b)\bigr).
\]
The parenthesis is fully symmetric. Every symmetric
homogeneous polynomial \(P\) occurs by choosing \(Q=P/3\),
so the map is onto. A basis of the symmetric space is
\(e_1^i e_2^j e_3^k\) with \(i+2j+3k=d-2\).
Counting these triples gives the stated dimension.
\end{proof}

\paragraph{Why the diagonal ceases to determine cyclic fifth derivatives.}
The preceding dimensions equal one at derivative orders three
and four, which explains why the source's cubic cyclic formula
and \eqref{qtr:cyclicquartic} each need just the corresponding
diagonal derivative. At order five the dimension is two.
This is not merely a count: the allowed deformation
\[
 \Delta R_4(A,B,C)=\frac{AB}{3}
 \bigl((A+B+C)^2-3(AB+BC+CA)\bigr)
\]
vanishes on the diagonal and preserves the axis and Euler
relations, but changes the cyclic fifth derivative by
\[
 5!\,abc\bigl((a+b+c)^2-3(ab+bc+ca)\bigr).
\]
Therefore the diagonal fifth derivative together with the
named relations cannot determine every cyclic fifth ray.
Additional analytic information is needed.

\subsection{An exact symmetric generating germ}

The preceding cyclic statement can be expressed without
choosing coordinates for each degree.

\begin{proposition}[Cyclic factorization]\label{qtr:factorization}
There is a unique symmetric holomorphic germ \(J(A,B,C)\)
near the origin such that
\begin{align}
 \sum_{\mathrm{cyc}}T(A,B,C)
 ={}&2\bigl(\zeta(A)\zeta(B)+\zeta(B)\zeta(C)
                              +\zeta(C)\zeta(A)\bigr)\notag\\
 &-\zeta(A+B)-\zeta(B+C)-\zeta(C+A)\notag\\
 &+2\bigl(\zeta(A)+\zeta(B)+\zeta(C)\bigr)+1
 +ABC J(A,B,C).
 \label{qtr:J}
\end{align}
\end{proposition}

\begin{proof}
In the cyclic sum of \eqref{qtr:polar}, the two terms with
denominator \(A+C\) combine to \(-\zeta(B-1)\), and likewise
for the other two denominators. Thus the cyclic sum is
holomorphic and symmetric near the origin.
On \(B=0\), the Euler relation gives
\[
 T(A,0,C)+T(0,C,A)+T(C,A,0)
                   =2\zeta(A)\zeta(C)-\zeta(A+C).
\]
The explicit zeta expression in \eqref{qtr:J} has exactly the
same restriction. By symmetry, their difference vanishes on
all three coordinate planes. The convergent Taylor series is
therefore divisible by \(ABC\). Its quotient is holomorphic
and symmetric, and uniqueness is immediate.
\end{proof}

Let \(J_k\) denote the homogeneous degree-\(k\) Taylor
polynomial of \(J\). For every \(n\ge3\) and every direction
whose pairwise sums are nonzero, differentiating the restricted
identity gives the all-order formula
\begin{align}
 \sum_{\mathrm{cyc}}W_{a,b,c}^{(n)}(0)
 ={}&2\sum_{\{u,v\}\in\{\{a,b\},\{b,c\},\{c,a\}\}}
       \sum_{j=1}^{n-1}\binom nj u^jv^{n-j}z_jz_{n-j}\notag\\
 &-\bigl((a+b)^n+(b+c)^n+(c+a)^n\bigr)z_n
 +n!\,abc\,J_{n-3}(a,b,c).
 \label{qtr:allcyclic}
\end{align}
The inner sum is symmetric in \(u,v\), so the unordered-pair
notation is unambiguous. All derivatives at \(-1\) and all
boundary Gamma coordinates have disappeared. At the first
two orders of \(J\),
\[
 J_0=\frac{\Omega}{2}+3Lz_2+4z_3,\qquad
 J_1(A,B,C)=\frac{A+B+C}{24}
                  (\Lambda+8Lz_3-12z_2^2+16z_4),
\]
where \(\Omega=W_{1,1,1}'''(0)\) is the incoming cubic
diagonal coordinate. Thus \eqref{qtr:allcyclic} recovers the
source's cubic identity and proves
\eqref{qtr:cyclicquartic}.

\subsection{Verification, scope, and the next exact questions}

The accompanying \texttt{check\_quartic\_tornheim.py} verifies
\eqref{qtr:quarticformula}, \eqref{qtr:cyclicquartic},
the Euler specialization, and the elementary correction \(E\)
by exact symbolic arithmetic. It independently constructs the
homogeneous constraint matrices for degrees \(0\) through \(10\),
computes their exact ranks and cyclic images, and checks them
against the proved dimension formulas.
An optional numerical mode evaluates the new coordinates
through integrated Jonqui\`ere series and a Stirling-subtracted
Gamma sum, and compares asymmetric fourth derivatives with
an independent Mellin evaluation of \(T\). Numerical
discrepancies are diagnostics, not interval certificates.

The results answer the quartic directional-layer question and
the formal homogeneous-dimension part of the all-order
question in \cite{qtr:Independent}. They do not evaluate
\(\Lambda\), \(\Xi\), or the earlier cubic coordinates in
a finite standard-constant algebra, and they do not settle
\(S_6\) or the revised \(S_8\).
The Mellin and Jonqui\`ere ingredients are classical
\cite{DLMF}; the quartic and formal kernel results are
developments relative to the inspected repository snapshot.
A global priority claim is not made.

The next concrete problems are to identify a useful
independent analytic relation for the symmetric degree-two
germ \(J_2\); to construct an economical Gamma or Barnes
representation for \(\Xi\); and to compare additional
Tornheim functional relations with the kernel
\eqref{qtr:kernelform}. Any improved coordinate count should
specify exactly which extra relations are used.

% END sections/05_tornheim.tex


% BEGIN sections/06_audit.tex
\section{Source audit, corrections of interpretation, and integration}
\label{audit:section}

\subsection{What the audit establishes}

The targeted review found no false theorem in the selected incoming
Tornheim statements or in the harmonic continuation used here. It
did find an attribution overlap that matters for presenting this
contribution: the canonical manuscript already gives the polynomial
Hurwitz transform and the complete-homogeneous coefficient mechanism
for all Stieltjes moments. The precise locations are
\src{chapters/07-integration.tex},
\src{stieltjes:sec:moments},
\src{stieltjes:thm:transform}, and
\src{stieltjes:thm:moment} \cite{source:manuscript}.
Accordingly, Corollary~\ref{harm:Stieltjesmoments} is classified as a
fully indexed arbitrary-endpoint restatement and extension, useful
for the new harmonic reduction, rather than a new general principle.

The formal difference-equation input, the Espinosa--Moll generalized
polygamma, and the classical Barnes identity are also credited
explicitly. Likewise, the boundary Gamma representation for $U$ in
\eqref{qtr:U} is an existing theorem specialized to the derivative
order needed here. These distinctions prevent an exact consequence
from being counted twice as a separate research discovery.

The reviewed source files were taken from the fixed commit stated in
Section~\ref{scope:section}. As an additional provenance check, the
44 inspected TeX files from the three newest incoming reports were
compared byte-for-byte with the corresponding ZIP contents in that
checkout; every comparison agreed. The packaged manifest records
the nine archive names, their Git blob identifiers, and SHA-256
hashes. An integration based on a later repository revision should
reconcile later changes explicitly.

\subsection{Current conjectures must retain their current labels}

The canonical conjectures
\src{cycloquot:conj:S6} and \src{s8new:conj:S8}
remain conjectural. The previously rejected $S_8$ coefficient vector
is a different candidate from the revised one. A rejection of the
former is not a rejection of the latter, and a numerical residual
for either candidate is not an analytic proof.

Two different computational statements in the incoming collection
also need to remain separate. The enlarged modular search reported
in \emph{Independent Orders} was interrupted before exhausting its
matrix; that record alone establishes neither membership nor
nonmembership. The later \emph{Exact Resonant Identities} report
establishes an exact obstruction for the complete Cayley ideal
together with its specifically enumerated additional row family.
That is a valid statement about the stated relation space. It does
not exclude additional distribution, double-shuffle, or functional
relations, and it is not a nonvanishing theorem for the actual
period \cite{qtr:Independent,ghp:ExactResonant}.

The analogous qualification is essential for the present
Theorem~\ref{qtr:kernel}: the formal dimension is proved under
symmetry, the $C$-axis identity, and the Euler-plane identity.
Its exactness within that system does not imply arithmetic
independence, nor does it claim to have imposed every possible
Tornheim functional equation.

The earlier source corrections concerning Tin's equation (22),
local polylogarithm estimates, and the rejected old $S_8$ vector
were already reported before this contribution. They are not
advertised here as newly discovered errors.

\subsection{Specific normalization rules supported by the new proofs}

The following are concrete instructions for using the formulas.
They also identify the corrections required if one extends the
existing calculus by a naive parameter substitution.

\paragraph{Retain the entire resonant amplitude.}
For the resolvent family the local term is
$r_p(z)A_p(s)x^{s-1}$, with
$A_p(s)=\prod_{j=1}^p(s-j)$. Its continued integral is
$r_pA_p(s)/s$. Subtracting only its residue at zero leaves
incorrect regular coefficients. The correct expression is
\eqref{res:completion}, with the entire polynomial amplitude
retained. At $p=1$, failure to do so misses the explicit constant
$1$ in \eqref{res:explicitcheck}. This is a proved correction
to that tempting extension, not an allegation that the pinned
source already contains the incorrect resolvent formula.

\paragraph{Keep the endpoint coordinate.}
All endpoint finite parts in the article use $x$, with the
cutoff specified in Section~\ref{res:section}. A nonlinear
endpoint substitution can change the finite constant of a
logarithmic divergence. The resolvent's spectral answer is
converted back to this stated coordinate before coefficient
extraction. Cauchy boundary values in \eqref{res:jump} are
asserted for the ordinary $p=0$ integral; no unproved transport
of divergent higher-derivative finite parts is implicit.

\paragraph{Do not differentiate a fixed index as though it were free.}
The deletion theorem is proved on the hyperplane where a
specified index equals $-m$. It permits every derivative in
the remaining independent free indices, because it is an
entire identity in those indices. It does not determine the
normal derivative transverse to the eliminated-index
hyperplane. That additional information is a separate research
problem.

\paragraph{Keep both Bernoulli endpoint conventions.}
The strict-gap deletion rule uses $B_1(0)=-1/2$ and
$B_1(1)=1/2$. Identifying them incorrectly deletes the last
term in
\[
 \mathcal H(s,0,t)=\mathcal H(s-1,t)
                 -\mathcal H(s,t-1)-\mathcal H(s,t).
\]
The verification suite deliberately corrupts that sign and
confirms that an independent finite nested sum detects it.

\paragraph{Restrict individual Tornheim rays before specialization.}
The individual fourth derivative in \eqref{qtr:quarticformula}
requires $(a+c)(b+c)\ne0$. A sum of cyclic germs can extend
across planes where individual terms are singular. The
holomorphic cyclic identity \eqref{qtr:J} has that extension;
its representation as a sum of individually restricted
derivatives is used only when every pairwise direction sum
is nonzero.

\paragraph{Specify what is centered and what is simultaneous.}
The involution in Theorem~\ref{ghp:thm:classification} changes
all even harmonic tails simultaneously. It does not require
each tail to occur to an even power separately. For example
$(\zeta(2)-H_n^{(2)})(\zeta(4)-H_n^{(4)})$ qualifies.
Conversely, a raw even-order harmonic number does not qualify:
$H_n^{(2)}$ produces the pole $-1/s$ at $s=0$.
This prevents two opposite misreadings of the classification.

\subsection{Proposed integration into the research collection}

The material separates naturally into four independently
readable additions: the centered generalized-harmonic
classification and sums; rational cotangent transforms and
balanced primitives; arbitrary-depth harmonic deletion and
the Lerch generator; and quartic and higher Tornheim jets.
Each has distinct label prefixes in the supplied TeX.

The incoming source questions can be annotated with the exact
answers in Table~\ref{scope:table}. Q11 receives a complete
classification. Q8 receives the stated rational sector.
Q9 receives the balanced primitive ladder for these kernels,
while a closed formula for every unbalanced Barnes product
is still open. The quartic directional question receives a
full normal form, and the all-order question receives an
exact answer for its specified formal relation system.

No existing conjecture should be promoted to a theorem
unless it is exactly one of these proved statements.
The package includes a claim ledger, proposed placement
notes, source hashes, and separate exact and numerical
verification records so that this distinction can survive
later merging and renumbering.

% END sections/06_audit.tex


% BEGIN sections/07_research.tex
\section{Further research: precise next questions}
\label{research:section}

The proved results suggest several exact research problems.
They are listed with the additional information needed for
a solution. None is asserted on numerical evidence alone.

\subsection{Centered harmonic values and global spectral derivatives}

\paragraph{R1. Evaluate mixed even-tail moments in a controlled algebra.}
The invariant-ring theorem includes every product
\[
 (\zeta(2i)-H_n^{(2i)})(\zeta(2j)-H_n^{(2j)})
\]
and its products with odd-order harmonic numbers and powers
of $H_n$. Determine finite formulas for their centered
values at $s=-2r$, beginning with $i=1,j=2$.
The proof of Theorem~\ref{ghp:thm:square} suggests finite
summation by parts followed by Bernoulli coefficient
extraction. The target should specify an ordinary or
multiple-zeta coordinate space at each weight, with an
explicit account of the terms that survive. A useful
strong result would sum an entire family in one convergent
generator, as in \eqref{ghp:eq:polylog-generator}.

\paragraph{R2. Classify shifts other than one half.}
For a fixed positive rational $a\ne1/2$, determine all
polynomials $P$ whose principal parts vanish on the even
nonpositive spectral lattice. Proposition~\ref{ghp:prop:principal}
gives every finite obstruction, but the centered involution
no longer has its simple form. The key question is whether
the infinitely many Bernoulli constraints admit a finite
algebraic description. A classification of simultaneous
conditions at $a$ and $1-a$ may be more natural than a
classification at one noncentral shift.

\paragraph{R3. Separate local regularity from global first derivatives.}
The absence of a principal part is a local asymptotic
statement. The value of the first regular spectral
derivative also contains a global remainder integral.
Find polynomial combinations $P_1,\ldots,P_k$ and exact
coefficients for which these remainder kernels cancel
at a nonpositive even center. A satisfactory answer is
an identity of kernels or a normally convergent sum,
not a numerical relation among already evaluated
derivatives. The centered parity classification narrows
the possible local terms and should reduce this search.

\subsection{Beyond rational cotangent kernels}

\paragraph{R4. Complex powers with moving endpoint exponents.}
Extend the completed generator to
\[
 \operatorname{FP}_x\int_0^1
     \gamma_m^{(p)}(x)(K(x)-z)^{-\lambda}\,dx,
 \qquad z\notin\mathbb R.
\]
The local exponent is $\lambda-p-1$ before logarithmic
differentiation. Thus integer crossings of $\lambda$
create new resonances not present in a fixed repeated
resolvent. First choose a logarithm of $K(x)-z$
continuously along the interval, then prove a joint
meromorphic generator and identify each full resonant
amplitude before expanding either parameter. The rational
integer cases of Corollary~\ref{res:repeated} are exact
normalization data for this problem.

\paragraph{R5. Products of continuous polygamma orders.}
The transform \eqref{res:continuousformula} is linear in
the generalized function $\Psi(v,x)$. For products such
as $\Psi(v,x)\Psi(w,x)$, Fourier convolution introduces
ordered or Tornheim-type coordinates. Determine which
symmetric or reflected combinations reduce to ordinary
polylogarithmic order derivatives and which genuinely
retain a multivariate zeta coordinate. The useful
deliverable is an explicit joint generator together with
the contact terms on every resonant hyperplane.

\paragraph{R6. Unbalanced Barnes functions with fixed normalization.}
The balanced ladder gives combinations of
$\zeta'(-r,x)$ and $\zeta(-r,x)$, and hence normalized
Barnes combinations at the first levels. Seek individual
resolvent transforms of $\log G(x)$ and higher multiple
Gamma functions, with every polynomial-in-$x$ companion
term and normalization constant stated. A functional
equation alone does not determine the function without
the customary asymptotic normalization. One concrete
starting point is to combine \eqref{res:balancedmoment}
with transforms containing $x\log\Gamma(x)$, whose
Fourier coefficients can be computed by one integration
by parts.

\subsection{Entire harmonic interpolation at more than one free order}

\paragraph{R7. A jointly convergent generator for several free orders.}
The canonical gap polynomial is already explicit for
$k$ surviving free indices. Sum arbitrary zero blocks
before, after, and between these indices into a useful
multivariate analytic kernel. The coefficient identity
alone is insufficient: prove a domain in all decoration
variables, the cancellation of simultaneous Lerch
branches, and normal convergence of independent spectral
derivatives. The one-free-order formula
\eqref{harm:zerogenerator} supplies a stringent marginal
specialization for any proposed answer.

\paragraph{R8. The first transverse derivative at a deleted index.}
Differentiate a word in an index that will subsequently
be specialized to $-m$. Its lattice summand acquires
$-x^m\log x$, so a Bernoulli polynomial no longer
represents the gap sum. Identify a finite family of
new gap functions, perhaps in normalized Hurwitz-zeta
derivatives, closed under the same concatenation and
inversion argument. The first target is a fully proved
formula for
\[
 \left.\partial_u\mathcal H_{a,b}(s,u,t)\right|_{u=0}
\]
at arbitrary independent $s,t$. This would extend the
present theorem in a genuinely transverse direction.

\paragraph{R9. Rational endpoints and exact character projections.}
Combine the explicit endpoint coefficient formulas with
finite character decompositions at rational $a,b$.
The expected output is an algorithm expressing the
specified harmonic jets and polynomial Stieltjes moments
in Dirichlet $L$-derivatives and conductor logarithms.
The finite reduction should identify the character
parity, the normalization of the endpoint difference,
and every logarithmic derivative of the conductor
factor. This is particularly suited to exact symbolic
verification.

\subsection{The next Tornheim coordinates and the Gaussian conjectures}

\paragraph{R10. Determine the second-degree symmetric germ.}
The polynomial $J_2$ in \eqref{qtr:J} has two independent
formal coordinates under the named relations.
The diagonal fifth derivative provides only one linear
combination. Obtain one additional independent analytic
identity or a normalized ordinary convergent integral
that determines the other combination. The explicit
deformation following Corollary~\ref{qtr:cyclicimage}
proves why manipulating the existing diagonal and
Euler-plane relations alone cannot finish this task.

\paragraph{R11. Compare additional functional relations with the exact kernel.}
The representation \eqref{qtr:kernelform} is an explicit
target against which another Tornheim relation can be
tested. Expand a proposed functional relation at the
origin, compute its action on the $S$ and $Q$ components,
and state exactly which directions it removes. This
would turn a collection of individual ray identities
into a systematic all-order reduction theory. Any
improved dimension statement must list the added
relations; numerical dependence among sampled
coordinates is not a substitute.

\paragraph{R12. Give an independent Gamma or Barnes representation for \(\Xi\).}
The ordinary integral \eqref{qtr:Xi} is explicit, but it
does not yet explain whether the mixed coordinate
belongs to a smaller familiar constant algebra.
A useful next theorem would derive a different Gamma
or Barnes representation and prove its equality with
the polylogarithmic integral. Such an independent
representation may reveal functional symmetries or
exact cancellations that are hidden in either form.
Simply changing the splitting point in the same
subtracted integral does not provide new information.

\paragraph{R13. Find a functional bridge to the actual \(S_6\) and revised \(S_8\) candidates.}
The short Gaussian reductions remain important open
problems in the source. The completed bounded relation
obstruction shows that a proof must use information
beyond the particular imposed Cayley and extra-row
system. A constructive next step is an explicit
one-parameter identity whose Taylor coefficient is the
target mixed harmonic sum and whose second evaluation
uses the proposed Gaussian polylogarithmic constants.
All boundary terms, algebraic arguments, and logarithm
branches must be fixed before specialization. The
rational and Lerch generators in this article offer
possible tools, but no such bridge to either named
candidate has been proved here.

\subsection{A practical order of attack}

The most immediate exact projects are R1, R7, and R8:
each starts with a finite normal form or a convergent
kernel already supplied here. R10 and R11 are the
natural next steps for the directional-zeta program,
because the obstruction identifies precisely what
new information is missing. R4 and R5 require a more
substantial joint resonance analysis.

For formal verification, the first useful units are
the Bernoulli deletion identities, the canonical gap
normal form, the finite principal-part coefficients,
the quartic coefficient map, and the kernel decomposition
\eqref{qtr:kernelform}. These are exact finite algebraic
statements. Their analytic interfaces should then
be formalized separately: normal convergence,
endpoint subtraction, and the specified continuation
theorems. This separation makes a future machine-checked
development concrete without pretending that the
present numerical diagnostics already supply one.

% END sections/07_research.tex


% BEGIN sections/08_verification.tex
\section{Verification and reproducibility}
\label{verify:section}

\subsection{Analytic proofs, finite algebra, and numerical evidence}

The proofs in the article establish the stated infinite families.
Finite symbolic tests check the coefficient algebra and the
implementation of the formulas. Numerical tests compare independent
representations and help detect phase, endpoint, or transcription
errors. A small floating-point discrepancy is not an interval
enclosure and is not used as a proof of an identity.

The complete numerical records are supplied as JSON rather than
rounded tables alone. They record working precision, truncations,
directions or endpoints, and the actual observed discrepancies.
For a reproducible source baseline, the package also includes the
archive and canonical-file hashes and a claim ledger distinguishing
new developments, inherited inputs, explicit restatements, and
unresolved questions.

\subsection{Centered harmonic and trigamma-square checks}

The file \src{code/verify_harmonic_parity.py} verifies representative
invariant and anti-invariant polynomials through inverse degree 16,
the general-shift expansions at $a=1/3$ and $a=7/5$, and the formal
trigamma shift equation through the stated truncation. It independently
checks the two finite-summation constant terms for $j=1,\ldots,8$.
These finite tests supplement the all-order injectivity proof;
they do not replace it.

At 110-digit working precision, the values $\mathcal T(-2r)$ for
$r=0,\ldots,6$ were compared with a direct finite trigamma sum plus
an Euler--Maclaurin tail. The two cutoff choices were
$(N,J)=(64,38)$ and $(96,45)$. The largest discrepancy of the
tighter-cutoff value from the theorem was approximately
$1.51\cdot10^{-85}$; the largest difference between the two
cutoff calculations was $9.28\cdot10^{-69}$.
At $t=0.2,0.7,1.1$, the polylogarithmic generator was also compared
with the Bernoulli moment series and the ordinary integral generator.
The largest moment-series discrepancy was below $4.46\cdot10^{-95}$.
No rigorous remainder enclosure is inferred from these comparisons.

\subsection{Resolvent and continuous-order checks}

The file \src{code/verify_resolvent.py} performs 55 exact assertions
on rational local coefficients, the Gamma contact polynomial, its
harmonic derivative, and the threshold where that correction
vanishes. The completed numerical run has 28 checks at 55-digit
working precision. They include upper- and lower-half-plane
parameters, twelve pointwise polygamma finite parts, the explicit
trigamma correction, the first three balanced primitives at two
parameters, two convergent $\gamma_1'$ squared-resolvent integrals,
four direct Stieltjes integrals, and two continuous polygamma orders.
A special-value regression is also included.

For the ordinary and local-subtraction checks assigned tolerance
$10^{-42}$, the largest observed discrepancy was below
$3.31\cdot10^{-53}$. The endpoint-sensitive complex order
$v=0.3+0.2i$ had discrepancy approximately $8.04\cdot10^{-41}$,
within its separately stated $10^{-34}$ tolerance.
An additional independently implemented integration-by-parts
check in \src{code/verify_resolvent_contact.py} tests nine
polygamma contact values at 100 digits, with largest discrepancy
below $5.27\cdot10^{-97}$.

The initial numerical implementation used finite differences of
the Hurwitz function at spectral order zero. At the quadrature
node $x=7/8$, that evaluation lost precision even though the
mathematical formula was exact. The final program uses the
log-Gamma identity at this order and native Hurwitz order
derivatives with an exact endpoint shift elsewhere. The regression
at $7/8$ is retained, and no tolerance was relaxed. This is an
implementation correction, not a change to the theorem.

\subsection{Harmonic elimination, Lerch continuation, and primitives}

The exact reducer \src{code/reduce_harmonic.py} is exercised by
\src{code/verify_harmonic.py}. Its recorded checks comprise
26 words, 78 comparisons among different local deletion orders
and the canonical gap form, 52 independent finite nested sums,
327 coefficient degree conditions, seven zero-insertion
identities, and 36 complete-harmonic coefficient identities.
A deliberately corrupted strict-gap sign is detected.

At 60-digit precision, six arbitrary noninteger-endpoint
continuations and one free-order derivative were evaluated by
an independent Euler--Maclaurin procedure. Their largest
observed residual was below $1.43\cdot10^{-33}$.
The zero-decoration generating function, truncated at total
degrees $10,14,18$, was compared with separate Lerch evaluations;
the residuals were approximately
$3.53\cdot10^{-20}$, $1.53\cdot10^{-26}$, and
$7.82\cdot10^{-33}$, respectively.

An ordinary Mellin integral checks the Lerch formula in its
honest convergence domain. Three normalized polynomial
primitives and five Stieltjes quadratures give additional
independent checks, with discrepancies at approximately
59--62 decimal places. A Python reciprocal initially formed
at machine precision for an exact integer input was replaced
by arbitrary-precision division; the two affected cases were
rerun and the actual corrected records retained.

\subsection{Tornheim coefficient and Mellin checks}

The file \src{code/check_quartic_tornheim.py} checks the full
quartic coefficient formula, its cyclic form, diagonal and
Euler specializations, and the elementary correction in
the mixed coordinate $\Xi$ by exact symbolic arithmetic.
It independently constructs the symmetry, axis, and
Euler-plane constraint matrices for remainder degrees
$0$ through $10$, computes their ranks and cyclic images,
and compares them with the dimension formulas.

The numerical mode uses 70-digit working precision, 68
integrated Jonqui\`ere terms, an incomplete-Gamma cutoff
of 165, and a 16-node fourth-derivative formula with step
$0.000015$. The independent Mellin comparisons were:
\begin{center}
\begin{tabular}{@{}cc@{}}
\toprule
Direction $(a,b,c)$ & Absolute discrepancy\\
\midrule
$(1,2,3)$  & $3.739835213\cdot10^{-37}$\\
$(2,3,1)$  & $2.075573689\cdot10^{-38}$\\
$(1,-2,3)$ & $2.871059279\cdot10^{-39}$\\
\bottomrule
\end{tabular}
\end{center}
Reducing the finite-difference step by a factor of ten
from an initial run decreased the first discrepancy by
the expected scale for the retained truncation error.
The delivered record is the completed tighter-step run.
The exact proof does not depend on that numerical tuning.

\subsection{How to rebuild and inspect the package}

The standalone file
\src{Harmonic_Parity_and_Resolvent_Identities.tex}
contains the complete article, including its bibliography.
It requires only a standard LaTeX installation with the
listed packages and can be compiled without network access.
The modular \src{article.tex}, \src{preamble.tex}, and
\src{sections/} files are also retained for integration.

Run \src{python3 build.py} to regenerate the standalone
source from the modules and compile the PDF.
Run \src{python3 code/verify_all.py} for the routine
exact and moderate numerical checks, or add \src{--full}
to rerun the slower independent Stieltjes, Lerch, and
Tornheim comparisons. The individual scripts remain
usable separately. Python dependencies and versions
are recorded in the package.

The PDF was compiled and inspected for unresolved references,
unavailable glyphs, overfull displays, and page layout.
The independent mathematical review notes are supplied
as review artifacts. They are an additional check by
separate derivations, not external peer review or a
machine-checked proof.

% END sections/08_verification.tex

\clearpage
\phantomsection
\addcontentsline{toc}{section}{References}
\begingroup
\small
\begin{thebibliography}{99}

% BEGIN references.tex
\bibitem{source:manuscript}
Vladimir Reshetnikov, \emph{ProveIt: Polylogarithms manuscript},
repository revision
\src{c2cb285b64273e4cb2c8d7a03c68a4a99ac812d4}, 11 October 2026.
In particular, the chapters on integration and Stieltjes moments,
harmonic Gamma coefficients, and the current $S_6$ and $S_8$ candidates.
\url{https://github.com/VladimirReshetnikov/ProveIt/tree/c2cb285b64273e4cb2c8d7a03c68a4a99ac812d4/Analysis/Polylogarithms/docs/manuscript}.

\bibitem{ghp:ExactResonant}
ProveIt research collection, \emph{Exact Resonant Identities},
incoming package \src{ProveIt_Exact_Resonant_Identities_2026-10-11.zip},
especially the reflected-moment and raw harmonic-power sections,
the source audit, and research questions Q8, Q9, and Q11.
Inspected at the revision in \cite{source:manuscript}.
Git blob \src{715a1fd8e26469c5a02d91f8fec1f65a1848fde4}.
\url{https://github.com/VladimirReshetnikov/ProveIt/tree/c2cb285b64273e4cb2c8d7a03c68a4a99ac812d4/docs/incoming}.

\bibitem{qtr:Independent}
ProveIt research collection, \emph{Independent Orders and Exact Reductions},
incoming package \src{ProveIt_Independent_Orders_2026-10-11.zip},
especially the entire harmonic interpolation, the Tornheim analytic
normalization and boundary Gamma series, and the quartic and higher-order
research questions. Inspected at the revision in \cite{source:manuscript}.
Git blob \src{a8b455bd3e9139c6930a04da92c4c5a11b300393}.

\bibitem{source:Cubic}
ProveIt research collection, \emph{Cubic Tornheim Derivatives and Harmonic
Stieltjes Constants: A two-coordinate formula, a digamma-cube bridge,
and finite reductions at mixed integer orders},
incoming package \src{ProveIt_Cubic_Harmonic_Mixed_Orders_2026-10-11.zip}.
Inspected at the revision in \cite{source:manuscript}.
Git blob \src{428f844851073b4b6eced0a9fece07aa64e948f1}.

\bibitem{source:Mixed}
ProveIt research collection, \emph{Mixed Spectral Identities},
incoming package \src{ProveIt_Mixed_Spectral_Identities_2026-10-10.zip},
especially \src{sections/03_dougall.tex}, Theorem
\src{dg:square} on weighted squares of trigamma differences.
Inspected at the revision in \cite{source:manuscript}.
Git blob \src{ddf0132e2700e443f321a5dc807f9737c37184bc}.

\bibitem{DLMF}
NIST Digital Library of Mathematical Functions,
\S5.11 (Gamma asymptotics), \S5.15 (polygamma functions),
\S24.4 (Bernoulli polynomial identities),
\S25.11 (Hurwitz zeta), and \S25.12 (polylogarithms).
In particular, equations 25.11.12, 25.11.14, 25.11.17,
25.11.18, and 25.12.12. Accessed 11 October 2026.
\url{https://dlmf.nist.gov/5.11};
\url{https://dlmf.nist.gov/5.15};
\url{https://dlmf.nist.gov/24.4};
\url{https://dlmf.nist.gov/25.11};
\url{https://dlmf.nist.gov/25.12}.

\bibitem{ghp:ADH}
Boris Adamczewski, Thomas Dreyfus, and Charlotte Hardouin,
\emph{Hypertranscendence and linear difference equations},
Journal of the American Mathematical Society \textbf{34} (2021),
475--503. Theorem 1.2, case $\mathbf S_\infty$.
\url{https://doi.org/10.1090/jams/960};
\url{https://arxiv.org/abs/1910.01874}.

\bibitem{harm:INY}
Kentaro Ihara, Yayoi Nakamura, and Shuji Yamamoto,
\emph{Interpolant of truncated multiple zeta functions},
arXiv:2407.20509 (2024).
\url{https://arxiv.org/abs/2407.20509}.

\bibitem{EM}
Olivier Espinosa and Victor H. Moll,
\emph{A generalized polygamma function},
Integral Transforms and Special Functions \textbf{15} (2004),
101--115. The generalized order function and balanced
negapolygamma normalization.
\url{https://arxiv.org/abs/math/0305079}.

\bibitem{harm:EM}
Olivier Espinosa and Victor H. Moll,
\emph{On some integrals involving the Hurwitz zeta function: Part 2},
The Ramanujan Journal \textbf{6} (2002), 449--468.
Theorem 2.1 and Example 2.2.
\url{https://www.math.tulane.edu/~vhm/papers_html/hurwitz2.pdf}.

\bibitem{Vardi}
Ilan Vardi,
\emph{Determinants of Laplacians and multiple gamma functions},
SIAM Journal on Mathematical Analysis \textbf{19} (1988),
493--507. The Hurwitz-zeta derivative and Barnes-$G$ relation.
\url{https://doi.org/10.1137/0519035}.

% END references.tex

\end{thebibliography}
\endgroup
\end{document}
